% !TeX root = ../../Book2.tex
\chapter{Fourier 乘子}
\label{b2:ch:22}
\BookChapterNumber{b2:ch:22}
\begin{chapterabstract}
有界符号在 Fourier 侧定义了一个 \(L^2\) 算子，但同一个算子是否保持 \(L^p\)，还取决于频率之间如何共同作用。本章先建立乘子的代数、对偶、极限与对称性质，再把环带上的导数或 Sobolev 控制转为奇异积分估计。应用部分区分零阶变换、分数阶正则性、同指数椭圆估计和低频逆算子，说明符号大小、光滑性与核的可积性分别能够提供什么。
\end{chapterabstract}

\begin{learninggoals}\begin{itemize}
\item\label{b2:item:22-goal-1}从有界符号的 \(L^2\) 算子出发定义 \(L^p\) 乘子，准确使用复合、对偶、极限及尺度对称，并区分强收敛与算子范数收敛。
\item\label{b2:item:22-goal-2}把 Mihlin 导数条件转为归一化环带核估计，说明低、高频求和与 \(L^2\) 起点各自的作用。
\item\label{b2:item:22-goal-3}理解 Hörmander 环带 Sobolev 条件怎样给出积分差分控制，识别严格指数条件实际用于哪一步。
\item\label{b2:item:22-goal-4}在 Riesz 变换、Bessel 位势及椭圆方程中选择匹配的空间，区分同指数有界性、正则性提升和低频存在性问题。
\end{itemize}\end{learninggoals}

对一个固定频率作乘法很简单，困难在于不同频率重新叠加后，空间中的 \(L^p\) 大小会如何变化。Plancherel 使 \(p=2\) 的答案完全由符号的本性上确界决定；其他指数没有同样的等距性质。即使符号处处具有模一，或者支撑在一个有界集合中，也不能只凭这些信息得出全部 \(L^p\) 界。

\ProseParagraphBreak

\secref{b2:sec:2201}先说明所讨论的是哪一个算子。对于有界可测符号，\(L^2\) 定义不需要预先寻找普通函数核；向有限指数 \(L^p\) 延拓时，再核对共同定义域上的一致性。无穷指数不能通过 Schwartz 函数的范数稠密性处理，而应使用有界 \(L^1\) 算子的伴随。恒等算子、平移以及区间投影也提醒我们：乘子理论不局限于具有可积卷积核的算子。

\ProseParagraphBreak

主定理将\chapref{b2:ch:18}的奇异积分方法与\chapref{b2:ch:19}的频率分解接起来。在每个环带上缩放到固定大小，再估计逆 Fourier 变换，便可把符号条件转为核条件。Mihlin 版本使用有限阶逐点导数，Hörmander 版本使用环带上的 Sobolev 范数；后者保留证明需要的积分差分控制，而不要求先得到极限核的逐点梯度界。单个环带的可积核估计不能不经求和论证就用于无限多个尺度。

\ProseParagraphBreak

阅读应用时，应先分开观察零频和无穷频率。两个符号可以在高频具有相同阶数，却因零频行为不同而作用于不同的空间；对最高阶导数的先验界，也不自动提供一个同指数的无质量逆算子。本章所需预备知识是 Fourier 反演、Plancherel、Young 不等式、插值与对偶，以及已建立的 Calderón--Zygmund 和平方函数估计。章末题沿这些工具安排，不要求预先使用后章的球面乘子结论。

% !TeX root = ../../Book2.tex
\section{乘子算子}
\label{b2:sec:2201}

常系数微分算子经过 Fourier 变换成为多项式乘法，卷积算子则成为其核的 Fourier 变换所给出的乘法。求解方程时，还会遇到多项式的倒数、平方根或更一般的函数。因此，一个共同的问题是：对每个频率分别乘上 \(m(\xi)\)，所得操作能否在原来的函数空间中保持有界？

\ProseParagraphBreak
在 \(L^2\) 上，Plancherel 等式直接回答这个问题；在其他 \(L^p\) 上，频率之间的叠加不能由平方和完全描述，符号的振荡和不连续性便会产生影响。本节先把算子的定义和基本运算说清，再说明平移不变性为何引出乘子。下一节的尺度不变光滑性条件，将给出容易应用的一类充分判据。

\subsection{Fourier 乘子}
\label{b2:subsec:220101}

本章沿用
\[
 \widehat f(\xi)=\int_{\mathbb R^d}
 f(x)e^{-2\pi\mathrm i x\cdot\xi}\,dx
\]
的归一化。若 \(m\) 只是可测函数，一般不能把 \(m\widehat f\) 解释为任意两个分布的乘积。先取 \(m\in L^\infty(\mathbb R^d)\)，便可对每个 \(f\in L^2\) 定义
\[
 T_m f\coloneqq\mathcal F^{-1}(m\widehat f)\in L^2.
\]
这一步不要求 \(m\) 连续，也不要求它有普通函数意义下的逆 Fourier 变换。

\Needspace{6\baselineskip}
\begin{definition}[Fourier 乘子]\label{b2:def:22-multiplier}
设 \(d\geq1\)、\(1\leq p<\infty\)，且
\(m\in L^\infty(\mathbb R^d)\)。若存在常数 \(C\)，使以上在
\(L^2\) 中定义的算子对全部 \(f\in\mathcal S(\mathbb R^d)\) 满足
\[
 \|T_m f\|_p\leq C\|f\|_p,
\]
称 \(m\) 为 \(L^p\) \term[fourier-chengzi]{Fourier 乘子}[Fourier multiplier]，记 \(m\in M_p\)。该算子唯一延拓到
\(L^p\)，其范数记为 \(\|m\|_{M_p}\)。
对于 \(p=\infty\)，本章约定
\[
 m\in M_\infty\quad\Longleftrightarrow\quad
 \overline m\in M_1,
\]
并把 \(T_m\) 定义为 \(T_{\overline m}:L^1\to L^1\) 在
\(\langle f,g\rangle=\int f\overline g\) 配对下的伴随；
记 \(\|m\|_{M_\infty}=\|\overline m\|_{M_1}\)。
\end{definition}

当 \(p<\infty\) 时，唯一性来自 Schwartz 函数在 \(L^p\) 中的稠密性。
若 \(f\in L^2\cap L^p\)，先截断再作光滑化，可以找到在这两个范数中同时趋近 \(f\) 的 Schwartz 函数。因此，延拓后的 \(T_m f\) 与原来的 \(L^2\) 定义相同：两种范数收敛各自给出依测度收敛的子列，极限必一致。以后在共同定义域内使用 Fourier 计算，依据的就是这一相容性。

\ProseParagraphBreak
\(L^\infty\) 的情形必须另作说明。Schwartz 函数并不在其中范数稠密，例如常数 \(1\) 到任意趋于零的连续函数的 \(L^\infty\) 距离都不小于 \(1\)。上述伴随约定选出与 \(L^1\) 配对相容的 \WeakStar{}连续算子；它在 \(L^\infty\cap L^2\) 上也与 \(L^2\) 定义一致。事实上，对 \(g\in L^1\cap L^2\)，Plancherel 等式给出
\[
 \langle T_m f,g\rangle
 =\langle f,T_{\overline m}g\rangle,
\]
这个配对确定了同一个函数。以后涉及端点 \(\infty\) 的乘子，均采用此约定，而不从 Schwartz 类上的范数估计声称得到了任意意义下的唯一延拓。

\ProseParagraphBreak
最直接的例子来自卷积。若 \(k\in L^1\)、\(m=\widehat k\)，则
\[
 T_m f=k*f,\quad
 \|m\|_{M_p}\leq\|k\|_1,\quad 1\leq p\leq\infty.
\]
这里先在 Schwartz 类上用 Fubini 验证 Fourier 恒等式，再应用 Young 不等式。对于有限复 Borel 测度 \(\mu\)，同样有
\[
 m(\xi)=\int e^{-2\pi\mathrm i y\cdot\xi}\,d\mu(y),
 \quad T_m f(x)=\int f(x-y)\,d\mu(y),
 \quad \|m\|_{M_p}\leq|\mu|(\mathbb R^d).
\]
因此，常数符号 \(1\) 的核为 \(\delta_0\)，指数符号
\(e^{-2\pi\mathrm i y\cdot\xi}\) 对应平移。可积核只是乘子的一部分；Hilbert 变换的核则要用主值分布解释。下一节从局部化符号重构核时，首先得到的是原点外的普通函数，原点处可能存在的分布项仍由 \(L^2\) 定义确定。

\subsection{\texorpdfstring{\(L^2\)}{L2} 理论}
\label{b2:subsec:220102}

微分算子的符号通常无界，仍可在 \(L^2\) 的稠密子空间上定义。为准确地区别无界算子与有界乘子，对任意几乎处处有限的可测 \(m\)，令
\[
 D_m\coloneqq\{f\in L^2:m\widehat f\in L^2\},
 \quad T_m f=\mathcal F^{-1}(m\widehat f),\quad f\in D_m.
\]
该定义域稠密：把任意 \(\widehat f\in L^2\) 乘上
\(\mathbf1_{\{|\xi|\leq n,\ |m(\xi)|\leq n\}}\)，再令 \(n\to\infty\)，即可由支配收敛得到 \(L^2\) 逼近。

\Needspace{6\baselineskip}
\begin{proposition}\label{b2:prop:22-l2}
设 \(m:\mathbb R^d\to\mathbb C\) 可测且几乎处处有限。在自然定义域
\(D_m\) 上定义的 \(T_m\) 能延拓为 \(L^2\) 有界算子，当且仅当
\(m\in L^\infty\)。此时
\[
 \|m\|_{M_2}=\|m\|_\infty.
\]
\end{proposition}
\begin{proof}
若 \(m\in L^\infty\)，Plancherel 等式直接给出
\[
 \|T_m f\|_2^2=\int|m(\xi)|^2\cdot|\widehat f(\xi)|^2\,d\xi
 \leq\|m\|_\infty^2\|f\|_2^2.
\]
反过来，设在 \(D_m\) 上有界，范数上界为 \(C\)。对任意有限测度集合 \(E\)，令
\[
 E_n=E\cap\{|\xi|\leq n,\ |m(\xi)|\leq n\}.
\]
取 \(\widehat f=\mathbf1_{E_n}\) 时确有 \(f\in D_m\)，故
\(\int_{E_n}|m|^2\leq C^2|E_n|\)。由单调收敛，
\(\int_E|m|^2\leq C^2|E|\)。
如果 \(\{|m|>C+\epsilon\}\) 有正测度，取其中有限正测度子集代入，就产生矛盾。
因此 \(\|m\|_\infty\leq C\)，与充分性合起来得到范数等式。
\end{proof}

例如 \(|m|=1\) 时，\(T_m\) 在 \(L^2\) 上等距；若 \(m=\mathbf1_E\)，则 \(T_m\) 是频率集合 \(E\) 上的正交投影。无论 \(E\) 的边界多么复杂，这个 \(L^2\) 结论都成立。其他 \(p\) 的困难不在逐点乘法本身，而在 Fourier 变换不能再把原空间范数等同于频率空间的 \(L^p\) 范数。

\ProseParagraphBreak
有界符号的最大值还是任何乘子范数的下界。证明的办法是把 Fourier 质量集中在一个频率附近：在空间侧用越来越缓慢变化的函数，再乘上一个固定频率的指数。此处不能直接代入 \(e^{2\pi\mathrm i x\cdot\xi_0}\)，因为它不属于有限指数的 \(L^p(\mathbb R^d)\)。

\Needspace{6\baselineskip}
\begin{proposition}[符号的必要界]\label{b2:prop:22-symbol-bound}
设 \(d\geq1\)、\(1\leq p\leq\infty\)，且 \(m\in M_p\)，其中
\(M_\infty\) 采用\cref{b2:def:22-multiplier}的伴随约定。则
\[
 \|m\|_\infty\leq\|m\|_{M_p}.
\]
\end{proposition}
\begin{proof}
先设 \(1<p<\infty\)，\(p'=p/(p-1)\)。取
\(f(x)=e^{-\pi|x|^2}\)、\(g(x)=f(x)^{p-1}\)，并在 \(m\) 的 Lebesgue 点
\(\xi_0\) 处定义
\[
 f_\epsilon(x)=\epsilon^{d/p}e^{2\pi\mathrm i x\cdot\xi_0}f(\epsilon x),
 \quad
 g_\epsilon(x)=\epsilon^{d/p'}e^{2\pi\mathrm i x\cdot\xi_0}g(\epsilon x).
\]
两者的相应范数不随 \(\epsilon\) 改变。变换变量和 Plancherel 配对给出
\[
 \langle T_m f_\epsilon,g_\epsilon\rangle
 =\int m(\xi_0+\epsilon\eta)\widehat f(\eta)
       \overline{\widehat g(\eta)}\,d\eta.
\]
两个 Fourier 变换都是正 Gaussian 函数；其乘积除以积分后是一个
Gaussian 近似恒等核。Lebesgue 点性质因此使右端趋于
\(m(\xi_0)\int f\overline g\)。另一方面，左端的绝对值不超过
\(\|m\|_{M_p}\cdot\|f\|_p\cdot\|g\|_{p'}\)。
所选 \(g\) 恰使
\(\int f\overline g=\|f\|_p\cdot\|g\|_{p'}\)，故
\(|m(\xi_0)|\leq\|m\|_{M_p}\)。

当 \(p=1\) 时，仍取同一个 \(f\)，改用
\(g_\delta(x)=e^{-\pi\delta|x|^2}\)，其中 \(\delta>0\)。
对每个固定 \(\delta\)，以上计算以及 \(\|g_\delta\|_\infty=1\) 给出
\[
 |m(\xi_0)|\int f g_\delta\leq\|m\|_{M_1}\cdot\|f\|_1.
\]
再令 \(\delta\downarrow0\)，由支配收敛得到所需不等式。
\(p=\infty\) 的结论应用于 \(\overline m\in M_1\) 即得。
\end{proof}

这里从 Lebesgue 点得到的是几乎处处的结论，恰好对应本质上确界。
某个频率点上的任意赋值不影响乘子。这个下界通常不是充分条件：
\(L^2\) 的特殊性正是只凭它就能完全刻画有界性。

\subsection{\texorpdfstring{\(L^p\)}{Lp} 乘子问题}
\label{b2:subsec:220103}

在判别具体符号以前，先整理三种常用操作。它们使有限频率截断的计算能够进入无穷尺度的结论，也说明何时可以把两个已经构造的算子复合。

\Needspace{6\baselineskip}
\begin{proposition}[乘子的运算、对偶与极限]\label{b2:prop:22-multiplier-calculus}
设 \(d\geq1\)，并采用\cref{b2:def:22-multiplier}中的乘子定义。
\begin{enumerate}
 \item 对 \(1\leq p\leq\infty\)，若 \(m_1,m_2\in M_p\)，则
 \(m_1+m_2,m_1m_2\in M_p\)，且
 \[
 T_{m_1m_2}=T_{m_1}T_{m_2},\quad
 \|m_1m_2\|_{M_p}\leq\|m_1\|_{M_p}\cdot\|m_2\|_{M_p}.
 \]
 加法满足相应的三角不等式。
 \item 若 \(1<p<\infty\)，则
 \(m\in M_p\) 当且仅当 \(\overline m\in M_{p'}\)，并且
 \(\|m\|_{M_p}=\|\overline m\|_{M_{p'}}\)。
 两个算子在复 \(L^p\) 对偶下互为伴随。
 \item 对固定 \(1\leq p\leq\infty\)，若
 \(m_n\in M_p\)、\(\sup_n\|m_n\|_{M_p}\leq C\)，且
 \(m_n\to m\) 几乎处处，则 \(m\in M_p\)，
 \(\|m\|_{M_p}\leq C\)。
\end{enumerate}
\end{proposition}
\begin{proof}
先取 \(p<\infty\)。对 \(f\in L^2\cap L^p\)，
\(T_{m_2}f\) 仍在这两个空间内，故 Fourier 乘法可以连续使用两次：
\[
 \widehat{T_{m_1}T_{m_2}f}=m_1m_2\widehat f.
\]
两种延拓的相容性及 \(L^2\cap L^p\) 的稠密性给出算子恒等式与范数界。
加法的论证相同。

对第二条，在 \(f,g\in\mathcal S\) 上先用 \(L^2\) 配对写成
\[
 \langle T_m f,g\rangle
 =\int m\widehat f\,\overline{\widehat g}
 =\langle f,T_{\overline m}g\rangle.
\]
若 \(T_m\) 在 \(L^p\) 上有界，则右端作为 \(f\) 的泛函范数不超过
\(\|m\|_{M_p}\cdot\|g\|_{p'}\)。由 \(L^p\) 对偶性，
\(T_{\overline m}g\) 属于 \(L^{p'}\) 并满足相应界；
交换 \(p,p'\) 后得到范数等式。端点 \(p=\infty\) 的第一条则由
\(L^1\) 算子的伴随得到，注意标量乘子可交换，所以取伴随时复合次序的反转不改变结果。

对第三条，\cref{b2:prop:22-symbol-bound}先给出
\(\|m_n\|_\infty\leq C\)，因而 \(|m|\leq C\) 几乎处处。
当 \(f\in\mathcal S\) 时，支配收敛和 Plancherel 等式给出
\[
 \|T_{m_n}f-T_mf\|_2^2
 =\int|m_n-m|^2\cdot|\widehat f|^2\longrightarrow0.
\]
选取几乎处处收敛的子列。对于 \(p<\infty\)，Fatou 引理给出
\(\|T_mf\|_p\leq C\|f\|_p\)，正好是乘子定义所需的估计。
对于 \(p=\infty\)，把已经证明的 \(p=1\) 情形应用于
\(\overline{m_n}\)，再取伴随即可。
\end{proof}

极限结论只保证有界性，不保证算子范数收敛；即使符号逐点收敛，也可能在不断远移的频率处保留很大的差异。若 \(m\) 实值，对偶性给出 \(p\) 与 \(p'\) 的相同乘子范数。再与 \(L^2\) 界插值，可以得到已知指数与 \(2\) 之间的全部指数；这不能自动推出 \(1\) 或 \(\infty\) 端点。

\ProseParagraphBreak
乘子范数还具有若干精确的对称性。令
\[
 D_rf(x)=f(rx),\quad M_\eta f(x)=e^{2\pi\mathrm i\eta\cdot x}f(x),
 \quad \tau_yf(x)=f(x-y),\quad r>0.
\]
直接计算 Fourier 变换可得
\[
 \begin{aligned}
 T_{m(r\,\cdot)}&=D_{1/r}T_mD_r,\\
 T_{m(\,\cdot-\eta)}&=M_\eta T_mM_{-\eta},\\
 T_{e^{-2\pi\mathrm i y\cdot\xi}m(\xi)}&=\tau_yT_m.
 \end{aligned}
\]
\(D_r\) 的范数为 \(r^{-d/p}\)，其中 \(d/\infty=0\)，所以两个缩放范数相消。
调制和平移都是等距同构；把恒等式反向使用，得到这些操作都保持
\(\|m\|_{M_p}\)。正交坐标变换也有相同性质。对于端点 \(\infty\)，上述算子均与 \(L^1\) 配对相容，恒等式可由伴随验证。
这解释了下一节为什么估计
\(|\xi|^{|\alpha|}\partial^\alpha m(\xi)\)：缩放频率时，导数产生的尺度因子恰被前面的幂抵消。

\ProseParagraphBreak
一个不光滑但确实有界的例子是一维区间投影。Hilbert 变换 \(H\) 的符号为
\(-\mathrm i\,\operatorname{sgn}\xi\)，所以
\[
 P_+\coloneqq\tfrac12(I+\mathrm i H)
\]
在零点以外的符号为 \(\mathbf1_{(0,\infty)}\)。由\chapref{b2:ch:18}的
\(L^p\) 有界性，\(P_+\) 对 \(1<p<\infty\) 有界。若 \(a<b\)，则
\[
 T_{\mathbf1_{(a,b)}}=
 M_aP_+M_{-a}-M_bP_+M_{-b},
 \quad
 \|\mathbf1_{(a,b)}\|_{M_p}\leq2\|P_+\|_{p\to p},
\]
与区间的位置和长度无关。符号光滑性因此是充分条件而非必要条件；
也不能把这个一维结论直接推广到高维球的特征函数。边界曲率带来的困难将在\chapref{b2:ch:24}说明。

\ProseParagraphBreak
最后说明乘子与平移不变性的准确关系。卷积与平移交换很直观；
下面的反向结论说明，在 \(L^2\) 上，除此以外没有另一类有界平移不变算子。

\Needspace{6\baselineskip}
\begin{proposition}\label{b2:prop:22-translation-multiplier}
设 \(d\geq1\)。若复线性算子
\(T:L^2(\mathbb R^d)\to L^2(\mathbb R^d)\) 有界，且
\(T\tau_y=\tau_yT\) 对全部 \(y\in\mathbb R^d\) 成立，
其中 \(\tau_yf(x)=f(x-y)\)，则存在在几乎处处意义下唯一的
\(m\in L^\infty(\mathbb R^d)\)，使
\[
 T=T_m,\quad \|T\|_{2\to2}=\|m\|_\infty.
\]
反过来，每个有界 \(L^2\) 乘子都与所有平移交换。
\end{proposition}
\begin{proof}
令 \(A=\mathcal FT\mathcal F^{-1}\)。平移变成乘法
\(e^{-2\pi\mathrm i y\cdot\xi}\)，故 \(A\) 与所有这些乘法算子交换。
对任意 \(\varphi\in C_c^\infty\)，Fourier 反演把 \(\varphi\) 写成这些指数函数的
\(L^1\) 加权积分。更具体地，对每个 \(h\in L^2\)，
\[
 \varphi h=\int_{\mathbb R^d}
  (\mathcal F^{-1}\varphi)(y)e^{-2\pi\mathrm i y\cdot\xi}h(\xi)\,dy
\]
是 \(L^2\) 中的强积分，被积函数的范数不超过
\(|(\mathcal F^{-1}\varphi)(y)|\cdot\|h\|_2\)。
有界算子可移入这个积分，因而 \(A\) 与乘法 \(M_\varphi\) 交换。

接着把光滑截断换成可测集投影。若 \(E\) 有界，将 \(\mathbf1_E\) 与非负光滑近似恒等核卷积，可得
\(0\leq\varphi_n\leq1\)、\(\varphi_n\in C_c^\infty\)，且
\(\varphi_n\to\mathbf1_E\) 几乎处处。对每个 \(h\in L^2\)，
支配收敛说明 \(M_{\varphi_n}h\to M_{\mathbf1_E}h\)。
故 \(A\) 也与 \(M_{\mathbf1_E}\) 交换。再令有界截断集合递增到任意可测集，就得到一般情形。

取递增的有限测度立方体 \(Q_n\) 覆盖全空间。
由交换投影，\(A\mathbf1_{Q_n}\) 支持在 \(Q_n\) 内；记它为 \(m_n\)。
若 \(n<k\)，则
\[
 \mathbf1_{Q_n}m_k
 =A(\mathbf1_{Q_n}\mathbf1_{Q_k})=m_n.
\]
删去可数个零测集后，这些函数可以拼成可测函数 \(m\)。
对任意可测 \(E\subset Q_n\)，又有
\[
 A\mathbf1_E=\mathbf1_E m,\quad
 \int_E|m|^2\leq\|A\|^2|E|.
\]
取 \(E=\{|m|>\|A\|+\epsilon\}\cap Q_n\) 可知
\(\|m\|_\infty\leq\|A\|\)。由有限测度支集简单函数在 \(L^2\) 中的稠密性，
\(A=M_m\)。范数等式来自\cref{b2:prop:22-l2}；若两个符号给出同一算子，
将其差乘以任意有限测度集的特征函数，即得二者几乎处处相等。
反向结论只需注意两个逐点乘法算子可交换。
\end{proof}

该证明使用了 \(T\) 的平方界，尤其使用 Fourier 变换在 \(L^2\) 上的酉性。
若某算子最初只在另一个 \(L^p\) 空间给定，不能未经说明便把它搬到
\(L^2\) 上应用此命题。实际应用通常从一个有界符号出发，先建立
\(L^2\) 算子，再通过核估计证明它在其他 \(L^p\) 上也有界；这正是下面采用的顺序。

% !TeX root = ../../Book2.tex
\section{Mihlin–Hörmander 定理}
\label{b2:sec:2202}

有界符号在 \(L^2\) 上的作用已经由 Plancherel 完全描述；要得到其他指数上的界，还须控制不同频率的输出如何在空间中叠加。本节用两种条件完成这一步。Mihlin 条件逐点控制符号导数，先得到核的大小与梯度估计；Hörmander 条件只控制每个环带上的 Sobolev 范数，得到的是核差分的积分估计。两条路线都以 \(L^2\) 有界性为起点，再使用\chapref{b2:ch:18}的分解与插值方法。

\subsection{符号导数}
\label{b2:subsec:220201}

设 \(d\geq1\)，\(N\) 为非负整数，\(m\) 是定义在 \(\mathbb R^d\setminus\{0\}\) 上的复值 \(C^N\) 函数。尺度不变的导数条件写成
\[
 |\xi|^{|\alpha|}\cdot|\partial^\alpha m(\xi)|\leq A,
 \quad 0\leq|\alpha|\leq N,\quad \xi\ne0.
\]
这里 \(\alpha=(\alpha_1,\ldots,\alpha_d)\) 是多重指标，\(|\alpha|=\alpha_1+\cdots+\alpha_d\)，\(\partial^\alpha=\partial_1^{\alpha_1}\cdots\partial_d^{\alpha_d}\)。零阶条件就是 \(|m|\leq A\)。若观察半径约为 \(2^j\) 的频率环带，把它缩放到单位尺度后，一阶导数带来 \(2^j\)，而假设中的 \(|\xi|^{-1}\) 恰好抵消这个因子；更高阶导数遵循同一规律。

\ProseParagraphBreak
例如 \(m(\xi)=\xi_\ell/|\xi|\) 是零次齐次的：\(m(t\xi)=m(\xi)\)，\(t>0\)。将此式对 \(\xi\) 求导可得
\[
 (\partial^\alpha m)(t\xi)
 =t^{-|\alpha|}\cdot(\partial^\alpha m)(\xi).
\]
每个导数在单位球面上连续，因此任意固定有限阶的尺度导数界都成立。这里常数可以随阶数增长，并不要求一个常数同时控制所有阶导数。

\ProseParagraphBreak
另一个例子是 \(m(\xi)=b(\log|\xi|)\)，其中 \(b\in C^N(\mathbb R;\mathbb C)\)，并且 \(b,b',\ldots,b^{(N)}\) 全部有界。链式法则把 \(\partial^\alpha m\) 写成若干项之和，每一项都是 \(|\xi|^{-|\alpha|}\)、单位方向的有界光滑函数和某个 \(b^{(r)}(\log|\xi|)\) 的乘积，所以也满足尺度导数界。这样的符号可以沿尺度持续振荡，在原点不必有极限。乘子由几乎处处的符号决定，原点的取值不影响算子；定理要求的是离开原点后的光滑性。

\ProseParagraphBreak
尺度因子也划定了本节的范围。对 \(\alpha>0\)，\(|\xi|^{-\alpha}\) 在零频无界，因而不能是 \(L^2\) 上的有界零阶乘子；它适合放到改变可积性或正则性的分数次积分理论中。另一方面，普通的“各阶导数有界”未必包含高频处所需的衰减，不能直接代替尺度导数条件。

\subsection{二进分解}
\label{b2:subsec:220202}

先恢复\chapref{b2:ch:19}所用的分割。取实值径向函数 \(\chi\in C_c^\infty(\mathbb R^d)\)，使 \(0\leq\chi\leq1\)，在 \(|\xi|\leq1\) 时等于 \(1\)，在 \(|\xi|\geq2\) 时等于 \(0\)，并沿半径不增。令
\[
 \varphi(\xi)=\chi(\xi)-\chi(2\xi).
\]
于是 \(\varphi\geq0\)，其支撑包含于 \(\{\,1/2\leq|\xi|\leq2\,\}\)。有限和的望远镜恒等式为
\[
 \sum_{j=-M}^{N}\varphi(2^{-j}\xi)
 =\chi(2^{-N}\xi)-\chi(2^{M+1}\xi).
\]
对每个 \(\xi\ne0\)，令 \(M,N\) 趋于无穷便得到 \(\sum_{j\in\mathbb Z}\varphi(2^{-j}\xi)=1\)。每个非零频率只落在有限个相邻环带中，而且任何有限指标集 \(J\subset\mathbb Z\) 都满足
\[
 0\leq\sum_{j\in J}\varphi(2^{-j}\xi)\leq1.
\]
最后这个不等式将在取算子极限时提供共同的 \(L^2\) 控制。

\ProseParagraphBreak
把符号局部化后，令
\[
 \begin{aligned}
  m_j(\xi)&=m(\xi)\cdot\varphi(2^{-j}\xi),\\
  b_j(\zeta)&=m(2^j\zeta)\cdot\varphi(\zeta),\\
  h_j&=\mathcal F^{-1}b_j,\quad
  k_j(x)=2^{jd}\cdot h_j(2^jx).
 \end{aligned}
\]
这里第二行把所有频率环带送到同一个固定环带；Fourier 变量换元给出 \(k_j=\mathcal F^{-1}m_j\)。这样就可以先证明一个不依赖尺度的估计，再把尺度因子放回去。

\Needspace{6\baselineskip}
\begin{lemma}[固定环带上的有限导数核估计]\label{b2:lem:22-annular-kernel}
设 \(d\geq1\)，\(N\) 为非负整数，\(a\in C_c^N(\mathbb R^d;\mathbb C)\) 的支撑包含于 \(\{\,1/2\leq|\xi|\leq2\,\}\)，且
\[
 \max_{|\alpha|\leq N}\|\partial^\alpha a\|_\infty\leq M.
\]
则 \(h=\mathcal F^{-1}a\) 是光滑函数，并满足
\[
 |h(x)|+|\nabla h(x)|
 \leq C_{d,N}\cdot M(1+|x|)^{-N},
 \quad x\in\mathbb R^d.
\]
这里空间变量的一阶导数不要求符号具有 \(N+1\) 阶导数。
\end{lemma}
\begin{proof}
由于 \(a\) 支撑有界，任何多项式乘以 \(a\) 都可积。因此可以在逆 Fourier 积分内对 \(x\) 任意多次求导；这说明 \(h\) 光滑。在 \(|x|\leq1\) 时，直接估计积分给出 \(|h(x)|\leq C_dM\)。对 \(|x|>1\)，选取满足 \(|x_\ell|\geq|x|/\sqrt d\) 的坐标，分部积分 \(N\) 次：
\[
 h(x)=(-2\pi\mathrm i x_\ell)^{-N}
 \int_{\mathbb R^d}e^{2\pi\mathrm i x\cdot\xi}
 \cdot\partial_\ell^Na(\xi)\,\mathrm d\xi.
\]
边界项为零，因为 \(a\) 及所需低阶导数都紧支撑。积分区域的体积固定，于是
\[
 |h(x)|\leq C_{d,N}\cdot M|x|^{-N}.
\]
结合 \(|x|\leq1\) 的估计得到第一项的结论。

对于 \(\partial_{x_r}h\)，其频率符号是 \(2\pi\mathrm i\xi_r\cdot a(\xi)\)。在固定环带上，Leibniz 公式表明这个符号截至 \(N\) 阶的导数仍由 \(C_{d,N}M\) 控制：落在坐标 \(\xi_r\) 上的导数最多一次，落在 \(a\) 上的导数不超过 \(N\) 次。将刚才的估计用于该符号，再对 \(r=1,\ldots,d\) 求和，即得梯度界。
\end{proof}

现在对归一化符号 \(b_j\) 使用这个引理。Leibniz 公式明确给出
\[
 \partial^\alpha b_j(\zeta)
 =\sum_{\beta\leq\alpha}\binom{\alpha}{\beta}
 2^{j|\beta|}\cdot(\partial^\beta m)(2^j\zeta)
 \cdot\partial^{\alpha-\beta}\varphi(\zeta).
\]
在 \(\varphi\) 的支撑上，\(|\zeta|\) 介于 \(1/2\) 与 \(2\) 之间，因而
\[
 2^{j|\beta|}\cdot|(\partial^\beta m)(2^j\zeta)|
 \leq A|\zeta|^{-|\beta|}
 \leq 2^{|\beta|}A.
\]
所以 \(b_j\) 的 \(C^N\) 界与 \(j\) 无关。引理给出
\[
 \begin{aligned}
 |k_j(x)|&\leq C A\,2^{jd}(1+2^j|x|)^{-N},\\
 |\nabla k_j(x)|&\leq C A\,2^{j(d+1)}(1+2^j|x|)^{-N}.
 \end{aligned}
\]
本节中分割函数 \(\chi,\varphi\) 一经选定便保持不变，与它们有关的常数并入 \(C\)。

\ProseParagraphBreak
在 \(x\ne0\) 处，选整数 \(j_0\)，使 \(2^{j_0}|x|\leq1<2^{j_0+1}|x|\)。低频部分不需要衰减因子，高频部分则使用它。若 \(N>d\)，两段大小估计分别为
\[
 \begin{aligned}
 \sum_{j\leq j_0}2^{jd}
 &\leq C_d\,2^{j_0d}\leq C_d|x|^{-d},\\
 |x|^{-N}\sum_{j>j_0}2^{j(d-N)}
 &\leq C_{d,N}|x|^{-N}\cdot2^{j_0(d-N)}
 \leq C_{d,N}|x|^{-d}.
 \end{aligned}
\]
梯度估计的计算完全由同样两个级数给出，只需把其中的 \(d\) 改成 \(d+1\)，故要求 \(N>d+1\)。于是
\[
 \sum_{j\in\mathbb Z}|k_j(x)|\leq C A|x|^{-d},
 \quad
 \sum_{j\in\mathbb Z}|\nabla k_j(x)|
 \leq C A|x|^{-d-1}.
\]
在任意离开原点的紧集上，低频尾可由 \(\sum2^{jd}\) 与 \(\sum2^{j(d+1)}\) 一致控制，高频尾可由相应的衰减级数一致控制。因此核及其一阶导数的级数都局部一致收敛，
\[
 K(x)=\sum_{j\in\mathbb Z}k_j(x),\quad x\ne0,
\]
是 \(\mathbb R^d\setminus\{0\}\) 上的 \(C^1\) 函数，并满足上述大小与梯度界。这个核在离开原点处有普通函数意义，将用于给定 \(L^2\) 算子的支集外表示；其原点附近的奇性仍需由算子估计处理。

\subsection{Mihlin 乘子定理}
\label{b2:subsec:220203}

核的梯度界要求本证明选取 \(N>d+1\)。为避免在陈述中反复保留整数参数，以下采用 \(N=d+2\)。这得到\term[mihlin-chengzi-dingli]{Mihlin 乘子定理}[Mihlin multiplier theorem]的一种方便版本。

\Needspace{6\baselineskip}
\begin{theorem}[Mihlin 乘子定理的有限导数版本]\label{b2:thm:22-mihlin}
设 \(d\geq1\)，\(m\in C^{d+2}(\mathbb R^d\setminus\{0\};\mathbb C)\)，且
\[
 \sup_{\xi\ne0}|\xi|^{|\alpha|}\cdot|\partial^\alpha m(\xi)|
 \leq A<\infty,\quad |\alpha|\leq d+2.
\]
在原点任意指定 \(m\) 的值，并由
\(\widehat{T_mf}=m\cdot\widehat f\) 定义 \(L^2(\mathbb R^d;\mathbb C)\) 上的乘子。则对每个 \(1<p<\infty\)，\(T_m\) 唯一延拓为有界算子 \(L^p\to L^p\)，并且
\[
 \|T_mf\|_p\leq C_{d,p}\cdot A\|f\|_p.
\]
不同指数的延拓在相应空间的交集上一致。它还具有弱型 \((1,1)\) 延拓，弱型常数不超过 \(C_dA\)。
\end{theorem}
\begin{proof}
若 \(A=0\)，则 \(m=0\) 几乎处处，结论直接成立。以下设 \(A>0\)。零阶条件和\cref{b2:prop:22-l2}给出 \(\|T_m\|_{2\to2}\leq A\)。按前一小节构造 \(k_j\) 与 \(K\)，得到
\[
 |K(z)|\leq C A|z|^{-d},
 \quad |\nabla K(z)|\leq C A|z|^{-d-1},
 \quad z\ne0.
\]
还须证明这个 \(K\) 确实表示已经由 Fourier 变换定义的算子。令
\[
 m^{(L)}(\xi)=m(\xi)\cdot
 \sum_{|j|\leq L}\varphi(2^{-j}\xi),
 \quad K_L=\sum_{|j|\leq L}k_j.
\]
每个 \(k_j\) 都属于 \(L^1\)，故 \(T_{m^{(L)}}f=K_L*f\) 在 \(L^2\) 中成立。有限分割和介于 \(0\) 与 \(1\) 之间，且在非零频率趋于 \(1\)，于是
\[
 \|T_{m^{(L)}}f-T_mf\|_2^2
 =\int_{\mathbb R^d}|m^{(L)}-m|^2\cdot|\widehat f|^2
 \,\mathrm d\xi\longrightarrow0.
\]
这里共同控制函数是 \(4A^2\cdot|\widehat f|^2\)，可由支配收敛处理。

取紧支撑的 \(f\in L^2\)，并取紧集 \(E\subset\mathbb R^d\setminus\operatorname{supp}f\)。此时 \(f\in L^1\)，而所有差 \(x-y\)，其中 \(x\in E\)、\(y\in\operatorname{supp}f\)，组成离开原点的紧集 \(D\)。因此
\[
 \sup_{x\in E}\left|
 (K_L*f)(x)-\int K(x-y)\cdot f(y)\,\mathrm dy
 \right|
 \leq\|f\|_1\cdot\sup_{z\in D}|K_L(z)-K(z)|
 \longrightarrow0.
\]
另一方面，\(L^2\) 收敛有一个几乎处处收敛的子列。比较两个极限，得到 \(T_mf(x)=\int K(x-y)\cdot f(y)\,\mathrm dy\) 对 \(E\) 中几乎处处的 \(x\) 成立。以可数个这样的紧集穷尽支撑补集，便确认支撑外表示；该积分在支撑外绝对收敛。

把两变量核记为 \(\mathcal K(x,y)=K(x-y)\)。当 \(|x-x'|\leq|x-y|/2\) 时，连接 \(x-y\) 与 \(x'-y\) 的线段离原点至少 \(|x-y|/2\)，故微积分基本定理给出
\[
 |\mathcal K(x,y)-\mathcal K(x',y)|
 \leq C A\,\frac{|x-x'|}{|x-y|^{d+1}}.
\]
改变 \(y\) 得到另一变量的同样估计。于是 \(\mathcal K\) 满足标准 Calderón--Zygmund 核的大小与双变量光滑性条件，\(T_m\) 又已有 \(L^2\) 界。将\cref{b2:thm:18-cz-weak}和\cref{b2:thm:18-cz-lp}应用于 \(A^{-1}T_m\)，其核常数与 \(L^2\) 常数只依赖维数，得到所述弱型界与强型界。后者先在 \(L^2\cap L^p\) 上成立，再由该子空间在 \(L^p\) 中的稠密性唯一延拓；不同 \(p\) 的延拓用同一截断逼近识别。
\end{proof}

支撑外核的识别不能省略。例如常数符号 \(m=1\) 给出恒等算子，其分布核集中在原点，支撑外核却是零。一般地，支撑外核无法独自决定一个乘法项。本证明从 \(L^2\) 乘子出发，再核对它的支撑外表示，所以这种集中在对角线上的部分仍由 \(L^2\) 算子保留。

\ProseParagraphBreak
定理通常不能补上强型 \((1,1)\) 或 \(L^\infty\to L^\infty\) 两个端点。\chapref{b2:ch:18}的 Hilbert 变换给出了同时满足尺度导数条件而端点失败的例子。另一方面，若符号是支撑于一个固定环带的 \(C_c^{d+1}\) 函数，\cref{b2:lem:22-annular-kernel}已经使其核属于 \(L^1\)，Young 不等式便给出全部 \(1\leq p\leq\infty\) 的界。无限尺度之所以不同，是因为每个环带核的 \(L^1\) 范数虽一致有界，却未必对环带指标可求和。

\subsection{Hörmander 条件}
\label{b2:subsec:220204}

上面的论证把符号导数转成核的逐点衰减。若只知道导数的平方积分，仍能通过 Plancherel 得到核的带权平方积分；只要权足够强，就能控制核尾部的 \(L^1\) 积分。这条路线允许用非整数阶的光滑性，并避免要求极限核存在逐点梯度界。

\ProseParagraphBreak
本小节的 \(H^s\) 范数作用在符号的频率变量上。为免混淆，把该变量写成 \(\zeta\)，其对偶变量写成 \(x\)。对 \(s\geq0\)，采用
\[
 \|q\|_{H^s_\zeta}^2
 =\int_{\mathbb R^d}(1+4\pi^2|x|^2)^s
 \cdot|\mathcal F_\zeta q(x)|^2\,\mathrm dx.
\]
这是 \(p=2\) 的 Sobolev 范数。权关于 \(x\) 对称，因此用逆 Fourier 变换替代 \(\mathcal F_\zeta q\) 得到同一范数。

\Needspace{6\baselineskip}
\begin{lemma}[Sobolev 范数中的光滑截断]\label{b2:lem:22-sobolev-cutoff}
设 \(d\geq1\)，\(s\geq0\)，\(\psi\in C_c^\infty(\mathbb R^d;\mathbb C)\)。则乘法 \(q\mapsto\psi\cdot q\) 是 \(H^s(\mathbb R^d;\mathbb C)\) 上的有界算子，且
\[
 \|\psi\cdot q\|_{H^s}\leq C_{d,s,\psi}\cdot\|q\|_{H^s}.
\]
\end{lemma}
\begin{proof}
令 \(w_s(x)=(1+|x|)^s\)，它与 \((1+4\pi^2|x|^2)^{s/2}\) 等价。三角不等式蕴含
\[
 w_s(x)\leq w_s(x-y)\cdot w_s(y).
\]
先取 \(q\) 为 Schwartz 函数。乘积的 Fourier 变换是卷积，所以
\[
 w_s(x)\cdot|\widehat{\psi\cdot q}(x)|
 \leq
 \bigl((w_s\cdot|\widehat\psi|)*(w_s\cdot|\widehat q|)\bigr)(x).
\]
由于 \(\widehat\psi\) 是 Schwartz 函数，\(\|w_s\cdot\widehat\psi\|_1<\infty\)。Young 不等式给出
\[
 \|w_s\cdot\widehat{\psi\cdot q}\|_2
 \leq\|w_s\cdot\widehat\psi\|_1
 \cdot\|w_s\cdot\widehat q\|_2,
\]
即所需估计。

一般 \(q\in H^s\) 可在 \(H^s\) 中用 Schwartz 函数 \(q_n\) 逼近：在 Fourier 一侧先截断到有界区域，再以紧支撑光滑函数逼近，权在每个有界区域上与常数等价，逆变换即得所需 \(q_n\)。上式使 \(\psi\cdot q_n\) 在 \(H^s\) 中为 Cauchy 列；同时 \(q_n\to q\) 于 \(L^2\)，而 \(\psi\) 有界，所以 \(\psi\cdot q_n\to\psi\cdot q\) 于 \(L^2\)。两个极限一致，从而估计延续到 \(q\)。
\end{proof}

还需要把核的积分差分条件接到弱型估计上。为便于以后复用，下面单独给出卷积核版本。这个结果只要求支撑外表示和核差分积分有界，不要求核本身满足逐点大小条件。

\Needspace{6\baselineskip}
\begin{lemma}[积分 Hörmander 条件与弱型估计]\label{b2:lem:22-integral-hormander-weak}
设 \(d\geq1\)，\(T\) 是复线性有界算子 \(L^2(\mathbb R^d)\to L^2(\mathbb R^d)\)，\(\|T\|_{2\to2}\leq M\)。设 Borel 函数 \(K\in L^1_{\mathrm{loc}}(\mathbb R^d\setminus\{0\})\) 具有下列支撑外表示：对每个紧支撑的 \(f\in L^2\)，
\[
 Tf(x)=\int_{\mathbb R^d}K(x-y)\cdot f(y)\,\mathrm dy
\]
在 \(\operatorname{supp}f\) 外几乎处处成立，且右侧在这些点绝对收敛。若
\[
 H\coloneqq\sup_{v\ne0}
 \int_{|x|>2|v|}|K(x-v)-K(x)|\,\mathrm dx<\infty,
\]
则对 \(f\in L^1\cap L^2\) 和 \(t>0\)，
\[
 \bigl|\{\,x:|Tf(x)|>t\,\}\bigr|
 \leq \frac{C_d(M+H)}{t}\cdot\|f\|_1.
\]
因此 \(T\) 唯一延拓为与原 \(L^2\) 算子相容的弱型 \((1,1)\) 算子。
\end{lemma}
\begin{proof}
记 \(B=M+H\)。若 \(B=0\)，则 \(T=0\) 于 \(L^2\)，结论成立。设 \(B>0\)，并对 \(f\in L^1\cap L^2\) 在高度 \(\alpha=t/B\) 使用\cref{b2:thm:18-cz-decomposition}。该结果适用于复值函数，给出
\[
 f=g+\sum_Qb_Q,\quad
 \|g\|_2^2\leq 2^d\alpha\cdot\|f\|_1,\quad
 \sum_Q|Q|\leq\alpha^{-1}\cdot\|f\|_1,
\]
其中所选二进立方体两两不交，\(b_Q\) 支撑于 \(Q\)，\(\int b_Q=0\)，并且
\[
 \sum_Q\|b_Q\|_1\leq2\|f\|_1.
\]
每个坏块属于 \(L^2\)，坏块级数在 \(L^2\) 中收敛。这个收敛性质随后使有限坏块估计可以传到算子作用下的极限。

把每个 \(Q\) 关于其中心 \(c_Q\) 放大 \(8\sqrt d\) 倍，所得立方体记为 \(Q^*\)，并令
\[
 \Omega^*=\bigcup_QQ^*,\quad G=\mathbb R^d\setminus\Omega^*.
\]
由体积估计，
\[
 |\Omega^*|\leq C_d\cdot\frac{B}{t}\cdot\|f\|_1.
\]
好部用 \(L^2\) 有界性和 Chebyshev 不等式处理：
\[
 \bigl|\{\,|Tg|>t/2\,\}\bigr|
 \leq\frac{4M^2}{t^2}\cdot\|g\|_2^2
 \leq C_d\cdot\frac{M^2}{Bt}\cdot\|f\|_1
 \leq C_d\cdot\frac{B}{t}\cdot\|f\|_1.
\]

坏部所需的抵消正是积分条件发挥作用的地方。对 \(x\notin Q^*\)，支撑外表示和 \(\int b_Q=0\) 给出
\[
 Tb_Q(x)=\int_Q
 [K(x-y)-K(x-c_Q)]\cdot b_Q(y)\,\mathrm dy
\]
几乎处处成立。如果 \(y\in Q\) 且 \(x\notin Q^*\)，则 \(|x-c_Q|>2|y-c_Q|\)。令 \(z=x-c_Q\)、\(v=y-c_Q\)，Tonelli 定理和积分假设因此给出
\[
 \begin{aligned}
 \int_{\mathbb R^d\setminus Q^*}|Tb_Q(x)|\,\mathrm dx
 &\leq\int_Q|b_Q(y)|\cdot
   \left(\int_{\mathbb R^d\setminus Q^*}
    |K(x-y)-K(x-c_Q)|\,\mathrm dx\right)\,\mathrm dy\\
 &\leq H\|b_Q\|_1.
 \end{aligned}
\]
当 \(y=c_Q\) 时差分为零，其余点由 \(|z|>2|v|\) 的积分控制；不需为该单点另设条件。

枚举坏块并令 \(b^{(n)}\) 为前 \(n\) 项之和。在共同区域 \(G\) 上，有限和满足
\[
 \int_G|Tb^{(n)}(x)|\,\mathrm dx
 \leq H\sum_Q\|b_Q\|_1\leq2H\|f\|_1.
\]
由 \(b^{(n)}\to b=f-g\) 于 \(L^2\)，有 \(Tb^{(n)}\to Tb\) 于 \(L^2\)。选取几乎处处收敛的子列，再用 Fatou 引理，得到 \(\int_G|Tb|\leq2H\|f\|_1\)。于是
\[
 \bigl|\{\,x\in G:|Tb(x)|>t/2\,\}\bigr|
 \leq\frac{4H}{t}\cdot\|f\|_1.
\]
把这一界、好部界和 \(|\Omega^*|\) 相加，得到 \(L^1\cap L^2\) 上的弱型估计。

最后说明延拓的含义。对 \(f\in L^1\)，选有界紧支撑 \(f_n\)，使 \(\|f_n-f\|_1\leq2^{-3n}\)。对差 \(f_{n+1}-f_n\) 使用刚证的估计，可得
\[
 \sum_{n=1}^{\infty}
 \bigl|\{\,|Tf_{n+1}-Tf_n|>2^{-n}\,\}\bigr|<\infty.
\]
几乎每点最终都避开这些集合，故 \(Tf_n\) 几乎处处收敛。以其极限定义 \(Tf\)，再对 \(|Tf_n-Tf_m|\) 的分布估计令 \(m\to\infty\)，得到
\[
 \bigl|\{\,|Tf_n-Tf|>t\,\}\bigr|
 \leq\frac{C_dB}{t}\cdot\|f_n-f\|_1.
\]
同样的极限过程把弱型界传给 \(f\)，并说明不同逼近所得极限相同。用两组逼近的和与数乘可验证线性。若 \(f\in L^1\cap L^2\)，可选同时在 \(L^1,L^2\) 中逼近的截断，所得极限又与原算子的 \(L^2\) 极限相同。这给出相容性与弱型延拓的唯一性。
\end{proof}

固定 \(\eta\in C_c^\infty(\mathbb R^d\setminus\{0\})\)，使它在闭环带 \(\{\,1/2\leq|\zeta|\leq2\,\}\) 的一个邻域等于 \(1\)。对可测有界符号 \(m\)，令
\[
 B_s(m)\coloneqq\sup_{j\in\mathbb Z}
 \|m(2^j\,\cdot)\cdot\eta\|_{H^s}.
\]
固定截断 \(\eta\) 留出了一点环带边界余量，使 \(\varphi\) 和坐标乘以 \(\varphi\) 都能作为光滑乘法作用于被测量的对象。以下是\term[hormander-chengzi-dingli]{Hörmander 乘子定理}[Hörmander multiplier theorem]的环带版本。

\Needspace{6\baselineskip}
\begin{theorem}[Hörmander 环带条件]\label{b2:thm:22-hormander}
设 \(d\geq1\)，\(s>d/2\)，\(m\in L^\infty(\mathbb R^d;\mathbb C)\)。固定 \(\eta\in C_c^\infty(\mathbb R^d\setminus\{0\})\)，使它在 \(\{\,1/2\leq|\xi|\leq2\,\}\) 的邻域等于 \(1\)，并假设
\[
 B_s(m)=\sup_{j\in\mathbb Z}
 \|m(2^j\,\cdot)\cdot\eta\|_{H^s(\mathbb R^d)}<\infty.
\]
则 \(L^2\) 乘子 \(T_m\) 对每个 \(1<p<\infty\) 唯一延拓为有界算子 \(L^p\to L^p\)，且
\[
 \|T_mf\|_p\leq C_{d,p,s,\eta}
 \bigl(\|m\|_\infty+B_s(m)\bigr)\cdot\|f\|_p.
\]
它还具有与 \(L^2\) 定义相容的弱型 \((1,1)\) 延拓，弱型常数不超过
\(C_{d,s,\eta}\bigl(\|m\|_\infty+B_s(m)\bigr)\)。
\end{theorem}
\begin{proof}
仍用本节固定的 \(\varphi\)，并记
\[
 q_j(\zeta)=m(2^j\zeta)\cdot\eta(\zeta),\quad
 b_j=\varphi\cdot q_j,\quad h_j=\mathcal F^{-1}b_j,\quad
 k_j(x)=2^{jd}\cdot h_j(2^jx).
\]
因为 \(\eta=1\) 于 \(\operatorname{supp}\varphi\) 的邻域，\(b_j=m(2^j\,\cdot)\cdot\varphi\)，与二进分解中的符号相同。\cref{b2:lem:22-sobolev-cutoff}分别用于截断 \(\varphi\) 和 \(2\pi\mathrm i\zeta_r\cdot\varphi(\zeta)\)，给出
\[
 \|b_j\|_{H^s}+\sum_{r=1}^{d}
 \|2\pi\mathrm i\zeta_r\cdot b_j\|_{H^s}
 \leq C_{d,s,\eta}\cdot B_s(m).
\]
Plancherel 与 Sobolev 范数的定义因此说明
\[
 \|(1+|x|)^s\cdot h_j\|_2
 +\|(1+|x|)^s\cdot\nabla h_j\|_2
 \leq C\cdot B_s(m).
\]
这些 \(h_j\) 是光滑函数：\(b_j\) 有界且紧支撑，任何多项式乘以 \(b_j\) 都属于 \(L^1\)，可以在逆变换中对空间变量求导。光滑性来自有限频带上的可积性与空间变量求导。

令 \(\gamma=s-d/2>0\)。由 Cauchy--Schwarz 不等式，
\[
 \|h_j\|_1+\|\nabla h_j\|_1
 \leq C\cdot B_s(m)\cdot
 \left(\int_{\mathbb R^d}(1+|x|)^{-2s}\,\mathrm dx\right)^{1/2}
 \leq C_{d,s,\eta}\cdot B_s(m).
\]
更精确地，对 \(R\geq1\)，极坐标给出
\[
 \int_{|x|>R}(1+|x|)^{-2s}\,\mathrm dx
 \leq C_d\int_R^\infty r^{d-1-2s}\,\mathrm dr
 =\frac{C_d}{2s-d}\cdot R^{d-2s},
\]
所以
\[
 \int_{|x|>R}|h_j(x)|\,\mathrm dx
 \leq C_{d,s,\eta}\cdot B_s(m)R^{-\gamma}.
\]
这一步把带权 \(L^2\) 信息变成了两个不同用途的 \(L^1\) 信息：梯度的全空间积分控制小位移，核的尾积分控制大位移。

取有限指标集 \(J\subset\mathbb Z\)，令
\[
 K_J=\sum_{j\in J}k_j,\quad
 m_J(\xi)=m(\xi)\cdot\sum_{j\in J}\varphi(2^{-j}\xi),
 \quad T_Jf=K_J*f.
\]
每个 \(k_j\in L^1\)，故 \(K_J\in L^1\)，卷积与 \(L^2\) 乘子 \(T_{m_J}\) 一致。分割的非负性保证 \(|m_J|\leq\|m\|_\infty\)，因而
\[
 \|T_J\|_{2\to2}\leq\|m\|_\infty
\]
与有限集 \(J\) 无关。下一步证明积分 Hörmander 常数也与 \(J\) 无关。

固定 \(v\ne0\)。若 \(2^j|v|\leq1\)，由
\[
 k_j(x-v)-k_j(x)
 =-\int_0^1v\cdot\nabla k_j(x-tv)\,\mathrm dt
\]
及平移不变性可得
\[
 \begin{aligned}
 \int_{|x|>2|v|}|k_j(x-v)-k_j(x)|\,\mathrm dx
 &\leq|v|\cdot\|\nabla k_j\|_1\\
 &=2^j|v|\cdot\|\nabla h_j\|_1
 \leq C\cdot B_s(m)\cdot2^j|v|.
 \end{aligned}
\]
若 \(2^j|v|>1\)，则在 \(|x|>2|v|\) 上同时有 \(|x|>|v|\) 与 \(|x-v|>|v|\)。分别在两个积分中换元，使用上面的尾积分界，得到
\[
 \begin{aligned}
 &\int_{|x|>2|v|}|k_j(x-v)-k_j(x)|\,\mathrm dx\\
 &\quad\leq
 \int_{|z|>2^j|v|}|h_j(z)|\,\mathrm dz+
 \int_{|z|>2^j|v|}|h_j(z)|\,\mathrm dz
 \leq C\cdot B_s(m)\cdot(2^j|v|)^{-\gamma}.
 \end{aligned}
\]
选择 \(j_0\) 满足 \(2^{j_0}|v|\leq1<2^{j_0+1}|v|\)。两部分级数满足
\[
 \sum_{j\leq j_0}2^j|v|\leq2,\quad
 \sum_{j>j_0}(2^j|v|)^{-\gamma}
 \leq\frac{1}{1-2^{-\gamma}}.
\]
对任意有限 \(J\) 取子和不会增大这些上界，所以
\[
 \sup_{v\ne0}\int_{|x|>2|v|}
 |K_J(x-v)-K_J(x)|\,\mathrm dx
 \leq C_{d,s,\eta}B_s(m).
\]

有限和核是光滑的 \(L^1\) 函数，满足所需支撑外表示。于是\cref{b2:lem:22-integral-hormander-weak}给出所有 \(T_J\) 的一致弱型界。其 \(L^2\) 伴随的核为
\[
 K_J^*(x)=\overline{K_J(-x)}.
\]
变量替换 \(x\mapsto-x\)、\(v\mapsto-v\) 表明 \(K_J^*\) 具有同一个积分条件，伴随的 \(L^2\) 范数也相同，所以 \(T_J^*\) 具有同样的弱型界。记 \(D=\|m\|_\infty+B_s(m)\)；若 \(D=0\) 则结论平凡，否则先对 \(D^{-1}T_J\) 使用\cref{b2:thm:17-marcinkiewicz}，由弱型 \((1,1)\) 和强型 \((2,2)\) 得到 \(1<p<2\) 的界。对 \(D^{-1}T_J^*\) 作相同插值，再由 \(p'\in(1,2)\) 的对偶性得到 \(2<p<\infty\) 的界。因此
\[
 \|T_Jf\|_p\leq C_{d,p,s,\eta}\cdot D\|f\|_p,
 \quad1<p<\infty,
\]
常数与 \(J\) 的大小和位置无关。

最后令 \(J=\{\,-L,\ldots,L\,\}\)。分割恒等式和 \(|m_J|\leq\|m\|_\infty\) 使 \(T_Jf\to T_mf\) 于 \(L^2\)。对 \(f\in L^p\cap L^2\)，选取几乎处处收敛的子列，再用 Fatou 引理，将统一 \(L^p\) 界传到 \(T_mf\)。在 \(L^1\cap L^2\) 上也可把弱型分布估计传到极限：先把阈值 \(t\) 降为任意 \(t'<t\)，使用几乎处处收敛和 Fatou，再令 \(t'\uparrow t\)。强型延拓由 \(L^p\cap L^2\) 的稠密性得到，弱型延拓由前一引理证明末尾的逼近步骤得到；所有这些延拓与原 \(L^2\) 算子相容。
\end{proof}

严格不等式 \(s>d/2\) 在两处承担同一任务：使逆权 \((1+|x|)^{-s}\) 属于 \(L^2\)，并使尾积分指数 \(\gamma=s-d/2\) 为正。在 \(s=d/2\) 时，当前尾积分计算出现对数发散，高频环带级数也不再衰减。这指出了本证明不能到达端点的原因；证明失效本身并不构成某个附加条件下端点结论的反例。

\ProseParagraphBreak
有限环带和的估计通过 \(L^2\) 极限传到原算子。与 Mihlin 证明相比，保留下来的是足以估计坏块的核差分积分，而不是逐点梯度。这正是 Sobolev 环带条件可以降低光滑性要求的原因。

\ProseParagraphBreak
还可用平移乘子看清“充分条件”的含义。若 \(a\in\mathbb R^d\)，则 \(m(\xi)=e^{2\pi\mathrm i a\cdot\xi}\) 给出 \(T_mf(x)=f(x+a)\)，它在每个 \(L^p\) 上等距。但当 \(a\ne0\) 时，至少一个一阶导数的模恒为非零常数，乘上 \(|\xi|\) 后无界，所以不满足 Mihlin 条件。Mihlin 条件的作用在于直接检验符号，并使所得控制在尺度分解下保持一致。

\ProseParagraphBreak
最后给出下一节所需的一族乘子版本。标量算子范数一致有界并不能单独推出先在位置上求平方和的估计；此处统一的 Mihlin 条件还提供统一核结构，因而可以调用\cref{b2:prop:18-hilbert-valued-cz}。

\Needspace{6\baselineskip}
\begin{proposition}[统一 Mihlin 条件下的平方和估计]\label{b2:prop:22-uniform-mihlin-square}
设 \(d\geq1\)，\(1<p<\infty\)，\(I\) 是有限或可数指标集。对每个 \(\nu\in I\)，设 \(m_\nu\in C^{d+2}(\mathbb R^d\setminus\{0\};\mathbb C)\)，并假设
\[
 \sup_{\nu\in I}\sup_{\xi\ne0}
 |\xi|^{|\alpha|}\cdot|\partial^\alpha m_\nu(\xi)|
 \leq A<\infty,\quad|\alpha|\leq d+2.
\]
以各自 \(L^2\) 定义的相容 \(L^p\) 延拓记 \(T_{m_\nu}\)。则对任意复值函数族 \((f_\nu)_{\nu\in I}\)，只要右侧有限，就有
\[
 \left\|\left(\sum_{\nu\in I}|T_{m_\nu}f_\nu|^2\right)^{1/2}\right\|_p
 \leq C_{d,p}\cdot A
 \left\|\left(\sum_{\nu\in I}|f_\nu|^2\right)^{1/2}\right\|_p.
\]
常数与指标集的大小无关。
\end{proposition}
\begin{proof}
先设 \(I=\{1,\ldots,n\}\)，并令 \(H=\mathbb C^n\)，赋予 Euclidean 范数。在 \(L^2(\mathbb R^d;H)\) 上定义对角算子
\[
 \mathbb T(f_1,\ldots,f_n)
 =(T_{m_1}f_1,\ldots,T_{m_n}f_n).
\]
逐坐标 Plancherel 给出
\[
 \|\mathbb Tf\|_{L^2(H)}^2
 =\sum_{\nu=1}^{n}\int|m_\nu(\xi)|^2
 \cdot|\widehat f_\nu(\xi)|^2\,\mathrm d\xi
 \leq A^2\cdot\|f\|_{L^2(H)}^2.
\]
对每个符号，用本节二进分解构造支撑外核 \(K_\nu\)。全部符号共享常数 \(A\)，因此核的大小、梯度和支撑外表示也共享同一组常数。令
\[
 \mathbb K(z)=\operatorname{diag}
 \bigl(K_1(z),\ldots,K_n(z)\bigr),\quad z\ne0.
\]
对角矩阵在 \(H\) 上的算子范数等于对角元模的最大值，故
\[
 \|\mathbb K(z)\|_{H\to H}
 =\max_{1\leq\nu\leq n}|K_\nu(z)|
 \leq C A|z|^{-d}.
\]
当 \(|h|\leq|z|/2\) 时，各坐标的梯度界还给出
\[
 \|\mathbb K(z+h)-\mathbb K(z)\|_{H\to H}
 \leq C A\,\frac{|h|}{|z|^{d+1}}.
\]
对两变量核 \(\mathbb K(x-y)\)，改变任一变量都满足相应差分估计。支撑外表示由有限个坐标的表示逐项组合；此处有限维 Bochner 积分就是坐标积分。

这样，\(\mathbb T\) 满足\cref{b2:prop:18-hilbert-valued-cz}的全部条件：有 \(L^2(H)\) 起点、算子范数意义的标准核估计，以及紧支撑输入的支撑外表示。该命题的常数不依赖 Hilbert 空间的维数。应用它并对 \(A\) 归一化，即得
\[
 \|\mathbb Tf\|_{L^p(H)}
 \leq C_{d,p}\cdot A\|f\|_{L^p(H)},
\]
这正是有限族的平方和估计。先对 \(L^2(H)\cap L^p(H)\) 输入应用，再由稠密性推广，得到陈述中的相容 \(L^p\) 算子。

若 \(I\) 可数，枚举其元素并对前 \(n\) 项使用已证估计。右侧的有限平方和由完整平方和控制；左侧随着 \(n\) 增大逐点单调增加。对其 \(p\) 次方使用单调收敛定理，得到可数族结论。
\end{proof}

这个向量估计同时控制一组频带上的乘子。将它与\chapref{b2:ch:19}的平方函数定理结合，就能把逐频带的估计重新合成为原函数的范数；下一节将据此识别 Bessel 势范数。

% !TeX root = ../../Book2.tex
\section{应用}
\label{b2:sec:2203}

应用乘子定理时，先判断符号阶数，再选取匹配的函数空间。零阶符号可保持 \(L^p\)，负阶符号在不同频率区间可能改善正则性或改变可积指数，正阶符号则需输入具有相应导数。

\subsection{Riesz 变换}
\label{b2:subsec:220301}

Riesz 变换把一个方向上的微分除以总微分尺度。沿用\chapref{b2:ch:18}的归一化，令
\[
 r_j(\xi)\coloneqq-\mathrm i\frac{\xi_j}{|\xi|},
 \quad R_j\coloneqq T_{r_j},
 \quad 1\leq j\leq d,\quad \xi\ne0.
\]
先在 \(L^2(\mathbb R^d;\mathbb C)\) 中定义这个算子，零频处的取值不影响它。符号 \(r_j\) 为零次齐次函数，每次微分使齐次次数降低一。因此
\[
 (\partial^\alpha r_j)(t\omega)
 =t^{-|\alpha|}(\partial^\alpha r_j)(\omega),
 \quad t>0,\quad |\omega|=1.
\]
单位球面紧，右侧的角向因子有界，故满足\cref{b2:thm:22-mihlin}要求的 \(d+2\) 阶条件。这给出全部 \(1<p<\infty\) 上的有界延拓，并与\chapref{b2:ch:18}由奇异核得到的算子一致：两者在共同的 \(L^2\cap L^p\) 输入上相同，再由稠密性得到一致性。奇异积分证明直接显示核的零平均与尾部，乘子证明则便于把 Riesz 变换同微分算子组合。

\Needspace{6\baselineskip}
\begin{proposition}[Riesz 变换与导数]\label{b2:prop:22-riesz-gradient}
设整数 \(d\geq1\)、\(1<p<\infty\)，\(R_j\) 是符号 \(-\mathrm i\xi_j/|\xi|\) 的 \(L^p(\mathbb R^d;\mathbb C)\) 乘子。则
\[
 -\sum_{j=1}^dR_j^2f=f\quad(f\in L^p).
\]
对 \(f\in\mathcal S(\mathbb R^d)\)，令 \((-\Delta)^{1/2}f\) 为符号 \(2\pi|\xi|\) 的 Fourier 乘子，并把 \(\|\nabla f\|_p\) 理解为欧氏向量长度的 \(L^p\) 范数。有
\[
 \partial_jf=-R_j(-\Delta)^{1/2}f,\quad
 (-\Delta)^{1/2}f=\sum_{j=1}^dR_j\partial_jf,\quad
 \|(-\Delta)^{1/2}f\|_p\asymp\|\nabla f\|_p.
\]
此外，对每个 \(1\leq j,k\leq d\)，
\[
 \partial_j\partial_k f=R_jR_k(-\Delta f),
 \quad
 \sum_{j,k=1}^d\|\partial_j\partial_k f\|_p
 \leq C_{d,p}\|\Delta f\|_p.
\]
\end{proposition}
\begin{proof}
符号恒等式
\[
 -\sum_j\left(-\mathrm i\frac{\xi_j}{|\xi|}\right)^2=1
 \quad(\xi\ne0)
\]
先由 Plancherel 给出 \(L^2\) 上的算子恒等式。有限次复合在 \(L^p\) 有界，因此先限制到 \(L^2\cap L^p\)，再逼近一般 \(L^p\) 输入，得到第一式。

对 Schwartz 输入，\(\sum_jR_j\partial_jf\) 属于 \(L^p\cap L^2\)，其 Fourier 符号为
\[
 \sum_j\left(-\mathrm i\frac{\xi_j}{|\xi|}\right)
                (2\pi\mathrm i\xi_j)=2\pi|\xi|.
\]
这同时证明 \((-\Delta)^{1/2}f\in L^p\)，并给出
\(\|(-\Delta)^{1/2}f\|_p\leq C\sum_j\|\partial_jf\|_p\leq C'\|\nabla f\|_p\)。
再将 \(r_j\) 与 \(2\pi|\xi|\) 相乘，得到
\(R_j(-\Delta)^{1/2}f=-\partial_jf\)，所以反向估计也由 Riesz 有界性得到。这里两处负号由本书的负号 Fourier 变换约定决定。

最后，
\[
 \left(-\mathrm i\frac{\xi_j}{|\xi|}\right)
 \left(-\mathrm i\frac{\xi_k}{|\xi|}\right)
       4\pi^2|\xi|^2=-4\pi^2\xi_j\xi_k.
\]
因此二阶导数恒等式成立，对有限个复合算子应用 \(L^p\) 界即可。
\end{proof}
在 \(p=2\) 时，Plancherel 进一步给出准确等式
\(\|(-\Delta)^{1/2}f\|_2^2=\sum_j\|\partial_jf\|_2^2\)；一般 \(p\) 上则是范数可比。这个结果允许从 \(\Delta f\) 恢复每个二阶导数，而不仅是它们的和。对于 \(d=1\)，\(R_1\) 就是 Hilbert 变换，恒等式成为 \(H^2=-I\) 与 \((-\Delta)^{1/2}f=Hf'\)。区间示性函数的 Hilbert 变换有边界对数奇性及不可积的远场尾部，故这里的强型范围不能顺手补上 \(p=1,\infty\)。

\ProseParagraphBreak
还可处理任意光滑零次齐次符号 \(m\)，因为单位球面上的有限阶导数有界。零次齐次性并不要求球面平均为零；常数符号的核有一个原点质量。只有在提取原点外的奇异核时，才需要辨认其球面抵消，不能把抵消条件错误地加到所有零阶符号上。

\subsection{分数幂}
\label{b2:subsec:220302}

分数阶导数要同时指定高频的阶数与零频的处理。先对 Schwartz 函数记
\[
 J^sf\coloneqq\mathcal F^{-1}
       \bigl[(1+4\pi^2|\xi|^2)^{s/2}\cdot\widehat f\bigr]
 \quad(s\in\mathbb R),
\]
\[
 (-\Delta)^{s/2}f\coloneqq
 \mathcal F^{-1}\bigl[(2\pi|\xi|)^s\cdot\widehat f\bigr]
 \quad(s>0).
\]
前者是不齐次分数阶算子，后者是齐次分数阶算子。前者的符号在全空间光滑，各阶导数至多多项式增长，因而乘法保持 \(\mathcal S\)，并由对偶作用定义在全部 \(\mathcal S'\) 上；\(J^sJ^t=J^{s+t}\)。后者在零频不一定光滑，不能未经说明就与任意温和分布相乘。当前定义的 Schwartz 输入使频率乘积局部可积，已经足以进行下面的计算。

\ProseParagraphBreak

当 \(s>0\)，\(J^{-s}\) 的符号满足 Mihlin 条件，故在 \(1<p<\infty\) 上保持 \(L^p\)。齐次逆 \((-\Delta)^{-s/2}\) 在零频发散，应在适当齐次空间中定义，或在 \(0<s<d\) 时以\chapref{b2:ch:19}的 Riesz 势解释。后者在 \(1<p<d/s\) 时把 \(L^p\) 送入 \(L^q\)、\(1/q=1/p-s/d\)，改变了指数；不能据此声称它在同一个 \(L^p\) 上有界。

\Needspace{6\baselineskip}
\begin{definition}[Bessel 位势空间]\label{b2:def:22-bessel-space}
设整数 \(d\geq1\)、\(s\in\mathbb R\)、\(1<p<\infty\)。令 \(J^s\) 为符号
\((1+4\pi^2|\xi|^2)^{s/2}\) 的温和分布乘子。满足 \(J^sf\) 可由一个 \(L^p(\mathbb R^d;\mathbb C)\) 函数表示的温和分布组成\term[bessel-weishi-kongjian]{Bessel 位势空间}[Bessel potential space]，记作
\[
 H^{s,p}(\mathbb R^d)\coloneqq
 \{\,f\in\mathcal S'(\mathbb R^d):J^sf\in L^p(\mathbb R^d)\,\},
 \quad
 \|f\|_{H^{s,p}}\coloneqq\|J^sf\|_p.
\]
\end{definition}
这里把函数与它所定义的温和分布视为同一对象。这个空间完备，而且 \(J^s:H^{s,p}\to L^p\) 是等距双射，逆为 \(J^{-s}\)。事实上，若 \(g_n\to g\) 于 \(L^p\)，则对任意测试函数 \(\psi\in\mathcal S\)，
\[
 |\langle J^{-s}(g_n-g),\psi\rangle|
 \leq\|g_n-g\|_p\cdot\|J^{-s}\psi\|_{p'}
 \longrightarrow0.
\]
测试函数右侧仍为 Schwartz 函数。于是 \(J^{-s}g_n\) 在 \(\mathcal S'\) 中趋于 \(J^{-s}g\)，同时在所定义范数中收敛。取 \(g_n\in\mathcal S\) 逼近 \(g\)，还说明 Schwartz 类在 \(H^{s,p}\) 中稠密。相同的符号恒等式给出
\[
 J^t:H^{s,p}\longrightarrow H^{s-t,p},
 \quad \|J^tf\|_{H^{s-t,p}}=\|f\|_{H^{s,p}}.
\]
负阶空间允许分布输入；并不是每个元素都预先具有普通函数的代表。

\Needspace{6\baselineskip}
\begin{proposition}[Bessel 范数的平方函数描述]\label{b2:prop:22-bessel-square}
设整数 \(d\geq1\)、\(s\in\mathbb R\)、\(1<p<\infty\)。取实径向
\(\chi\in C_c^\infty(\mathbb R^d)\)，满足 \(0\leq\chi\leq1\)，在
\(|\xi|\leq1\) 为一、\(|\xi|\geq2\) 为零，并沿半径单调不增。令
\[
 \varphi=\chi-\chi(2\cdot),\quad
 P_0=\chi(D),\quad P_j=\varphi(2^{-j}D)\quad(j\geq1).
\]
这里 \(m(D)f=\mathcal F^{-1}(m\cdot\widehat f)\)，高频块 \(P_j\) 就是齐次块 \(\Delta_j\)，但 \(P_0\neq\Delta_0\)。对每个 \(f\in\mathcal S(\mathbb R^d)\)，有
\[
 \|J^sf\|_p\asymp
 \left\|\left(|P_0f|^2+
              \sum_{j\geq1}2^{2js}|P_jf|^2\right)^{1/2}\right\|_p.
\]
常数只依赖 \(d,p,s,\chi\)。更一般地，温和分布 \(f\) 的右侧有限当且仅当 \(f\in H^{s,p}(\mathbb R^d)\)。因此与\cref{b2:def:19-inhomogeneous-square}采用同一分割时，
\(F^s_{p,2}(\mathbb R^d)=H^{s,p}(\mathbb R^d)\)，且范数等价。
\end{proposition}
\begin{proof}
记 \(w_s(\xi)=(1+4\pi^2|\xi|^2)^{s/2}\)。先证明分析方向。令 \(g=J^sf\)，则
\(2^{js}P_jf=U_jg\)，其中
\[
 u_j(\xi)=2^{js}\varphi(2^{-j}\xi)w_{-s}(\xi).
\]
缩放到固定环带，得到
\[
 u_j(2^j\eta)=
 \varphi(\eta)(2^{-2j}+4\pi^2|\eta|^2)^{-s/2}.
\]
在 \(1/2\leq|\eta|\leq2\) 上，括号的值有与 \(j\geq1\) 无关的正下界和有限上界。反复求导只产生括号的实数幂与 \(\eta\) 的多项式，所以所有预先指定阶数的导数一致有界。

令 \(k_j=\check u_j\)。固定环带上的积分分部使归一化逆变换及其一阶导数具有任意阶一致衰减。因此对任意整数 \(M>d+1\)，
\[
 |k_j(x)|\leq C_M2^{jd}(1+2^j|x|)^{-M},\quad
 |\nabla k_j(x)|\leq C_M2^{j(d+1)}(1+2^j|x|)^{-M}.
\]
这正是\cref{b2:lem:19-uniform-dyadic-kernels}求和时使用的两种界；该求和只用一致常数，不要求归一化核逐个完全相同。将 \(2^j|x|\leq1\) 与 \(>1\) 分开求几何级数，得到
\[
 \sum_{j\geq1}|k_j(x)|\leq C|x|^{-d},\quad
 \sum_{j\geq1}|\nabla k_j(x)|\leq C|x|^{-d-1}.
\]

对有限指标集 \(E\subset\{1,2,\ldots\}\)，把 \(U_Eg=(U_jg)_{j\in E}\) 看作从标量到 \(\mathbb C^E\) 的算子。每个频率只落入固定数目的环带，故
\[
 \|U_Eg\|_{L^2(\ell^2)}^2
 =\int\sum_{j\in E}|u_j(\xi)|^2\cdot|\widehat g(\xi)|^2\,\mathrm d\xi
 \leq C\|g\|_2^2.
\]
列向量核的算子范数为 \((\sum_{j\in E}|k_j|^2)^{1/2}\)，不超过上面的绝对和；梯度同理。沿线段积分梯度给出两个变量的标准核差分界。有限个 Schwartz 核又保证支集外的普通卷积表示。因此\cref{b2:prop:18-hilbert-valued-cz}给出与 \(|E|\) 无关的 \(L^p\to L^p(\ell^2)\) 界。让 \(E\) 增加并用单调收敛，得到全部高频分析界。低频 \(P_0J^{-s}\) 的符号光滑紧支撑，单独由 Young 不等式控制；两者合起来给出平方范数不超过 \(C\|J^sf\|_p\)。

反向需要合成，而不能把各块的标量估计相加。取实偶的
\(\widetilde\varphi\in C_c^\infty(\{\,1/4<|\xi|<4\,\})\)，在
\(\operatorname{supp}\varphi\) 的邻域等于一；另取实偶
\(\widetilde\chi\in C_c^\infty\)，在 \(\operatorname{supp}\chi\) 的邻域等于一。令 \(V_0,V_j\) 的符号分别为
\[
 v_0=w_s\widetilde\chi,\quad
 v_j(\xi)=2^{-js}w_s(\xi)\widetilde\varphi(2^{-j}\xi).
\]
对高频族，归一化符号为
\[
 v_j(2^j\eta)=
 (2^{-2j}+4\pi^2|\eta|^2)^{s/2}\widetilde\varphi(\eta).
\]
与前面完全相同的固定环带导数界、有限重叠与列向量核证明给出
\[
 \left\|\left(\sum_{j\in E}|V_jh|^2\right)^{1/2}\right\|_{p'}
 \leq C\|h\|_{p'}.
\]
所有 \(v_j\) 均为实值，故 \(V_j\) 在积分配对下自伴随。对任意有限族 \(b_j\in L^p\)，逐点 Cauchy--Schwarz 与 Hölder 因而给出
\[
 \begin{aligned}
 \left|\int\left(\sum_{j\in E}V_jb_j\right)\cdot\overline h\right|
 &\leq\int
       \left(\sum_{j\in E}|b_j|^2\right)^{1/2}
       \cdot\left(\sum_{j\in E}|V_jh|^2\right)^{1/2}\\
 &\leq C\left\|\left(\sum_{j\in E}|b_j|^2\right)^{1/2}\right\|_p
          \cdot\|h\|_{p'}.
 \end{aligned}
\]
先在 Schwartz 输入上核对配对，固定有限族的 \(L^p\) 连续性再推广到一般 \(b_j\)。对偶性得合成界；低频项 \(V_0b_0\) 再用其 Schwartz 核单独估计。于是
\[
 \left\|V_0b_0+\sum_{j\in E}V_jb_j\right\|_p
 \leq C\left\|\left(|b_0|^2+\sum_{j\in E}|b_j|^2\right)^{1/2}\right\|_p.
\]
常数仍独立于有限指标集。

由于辅助截断在相应频率支撑上等于一，有限重构 \(f_N=P_0f+\sum_{j=1}^NP_jf\) 满足
\[
 J^sf_N=V_0P_0f+\sum_{j=1}^NV_j(2^{js}P_jf).
\]
对 Schwartz \(f\)，非齐次重构在 Schwartz 拓扑收敛。用刚才的合成界并取极限，得到所求反向不等式。

最后处理任意温和分布。若 \(g=J^sf\in L^p\)，取 \(g_n\in\mathcal S\) 在 \(L^p\) 中逼近它，令 \(f_n=J^{-s}g_n\)。分析界说明加权块序列在 \(L^p(\ell^2)\) 中收敛；逐个块在 \(\mathcal S'\) 中的极限又是 \(P_0f,2^{js}P_jf\)，故这个向量极限就是所需的频率块，得到右侧有限及上界。

反之，若右侧有限，每个块都是 \(L^p\) 函数。合成界可用于这些块，且对 \(N>M\geq0\)，
\[
 \|J^s(f_N-f_M)\|_p
 \leq C\left\|\left(
       \sum_{j=M+1}^N2^{2js}|P_jf|^2\right)^{1/2}\right\|_p
 \longrightarrow0.
\]
平方尾逐点趋零，并由题设的 \(L^p\) 平方函数控制。支配收敛说明 \(J^sf_N\) 在 \(L^p\) 中有极限。另一方面，\cref{b2:eq:19-inhomogeneous-reconstruction}给出 \(f_N\to f\) 于 \(\mathcal S'\)，而 \(J^s\) 在该空间连续，所以同一序列的分布极限是 \(J^sf\)。两个极限一致，证明 \(J^sf\in L^p\) 及反向界。
\end{proof}
估计常数独立于有限频带的个数，因而可以让频带数趋于无穷。平方和必须位于空间 \(L^p\) 范数内部；\(\sum_j2^{2js}\|P_jf\|_p^2\) 一般是另一种量。

\ProseParagraphBreak

整数阶时，频率权重应当恢复通常的弱导数。下面给出一个明确的符号恒等式，既处理混合导数，也避免把“椭圆性”留作没有写出的分解。

\Needspace{6\baselineskip}
\begin{proposition}[整数阶 Bessel 空间与 Sobolev 空间]\label{b2:prop:22-integer-sobolev}
设整数 \(d,k\geq1\)、\(1<p<\infty\)。令 \(W^{k,p}(\mathbb R^d;\mathbb C)\) 为所有弱导数 \(\partial^\alpha f\)、\(|\alpha|\leq k\)，均属于 \(L^p(\mathbb R^d)\) 的函数空间。则
\[
 H^{k,p}(\mathbb R^d)=W^{k,p}(\mathbb R^d),
\]
并且
\[
 \|J^kf\|_p
 \asymp\sum_{|\alpha|\leq k}\|\partial^\alpha f\|_p
 \asymp\|f\|_p+\sum_{j=1}^d\|\partial_j^kf\|_p.
\]
最后一项有限的 \(L^p\) 函数，已经自动具有全部至多 \(k\) 阶的 \(L^p\) 弱导数。
\end{proposition}
\begin{proof}
先从 \(H^{k,p}\) 出发。符号
\[
 \frac{(2\pi\mathrm i\xi)^\alpha}
      {(1+4\pi^2|\xi|^2)^{k/2}},
 \quad |\alpha|\leq k,
\]
在零点附近光滑有界；在高频处，\(\beta\) 阶导数为
\(O(|\xi|^{|\alpha|-k-|\beta|})\)。因而它满足 Mihlin 条件。把它作用于 \(J^kf\) 就是分布导数 \(\partial^\alpha f\)，得到第一方向及所有弱导数的存在。

反向只假定 \(f,\partial_1^kf,\ldots,\partial_d^kf\in L^p\)，其中导数按分布定义。令
\[
 q_j(\xi)=(2\pi\mathrm i\xi_j)^k,\quad
 Q(\xi)=1+\sum_{j=1}^d|q_j(\xi)|^2,\quad
 w_k(\xi)=(1+4\pi^2|\xi|^2)^{k/2}.
\]
有限维范数比较给出 \(Q(\xi)\asymp(1+|\xi|)^{2k}\)。多项式求导及恒等式 \(Q\cdot Q^{-1}=1\) 的反复微分说明
\[
 |\partial^\beta Q^{-1}(\xi)|
 \leq C_\beta(1+|\xi|)^{-2k-|\beta|}.
\]
例如一次微分为 \(-Q^{-2}\partial_\ell Q\)；更高阶每项都是 \(Q^{-1}\) 的若干因子乘上若干个 \(Q\) 的导数，总导数阶数仍为 \(|\beta|\)，给出所列幂次。

定义
\[
 a_0=\frac{w_k}{Q},\quad
 a_j=\frac{\overline{q_j}\,w_k}{Q}\quad(1\leq j\leq d).
\]
由上述界和 Leibniz 法则，\(a_0\) 的 \(\beta\) 阶导数为
\(O((1+|\xi|)^{-k-|\beta|})\)，\(a_j\) 的相应导数为
\(O((1+|\xi|)^{-|\beta|})\)，故全是 Mihlin 符号。精确恒等式
\[
 w_k=a_0+\sum_{j=1}^d a_jq_j
\]
于是给出分布等式
\[
 J^kf=T_{a_0}f+\sum_{j=1}^dT_{a_j}\partial_j^kf.
\]
这些光滑符号及其各阶导数至多多项式增长，分布乘法合法；其 \(L^p\) 延拓与该分布作用一致，可对 Schwartz 逼近使用对偶配对验证。右侧属于 \(L^p\)，且范数至多为题中最后一项的常数倍。故 \(f\in H^{k,p}\)，前半段再恢复全部混合弱导数。
\end{proof}
例如 \(k=2\) 时，只知道 \(f,\partial_1^2f,\ldots,\partial_d^2f\in L^p\)，就能控制所有 \(\partial_j\partial_kf\)，也能控制一阶导数。低阶项 \(\|f\|_p\) 保留了零频信息。该结论在整个 \(\mathbb R^d\) 上利用 Fourier 乘子建立，不自动给出任意区域或任意边界条件下的同一结论。

\ProseParagraphBreak

分数阶算子还具有一个反映非局部性的实空间公式。对 \(0<s<2\)，设
\[
 A_{d,s}=\int_{\mathbb R^d}\frac{1-\cos(2\pi z_1)}{|z|^{d+s}}\,\mathrm dz,
 \quad c_{d,s}=\frac{(2\pi)^s}{A_{d,s}}.
\]
原点处 \(1-\cos(2\pi z_1)\leq C|z|^2\)，无穷处它有界，故 \(0<A_{d,s}<\infty\)；这同时解释范围 \(0<s<2\)。
\Needspace{6\baselineskip}
\begin{proposition}[分数阶 Laplace 算子的二阶差公式]\label{b2:prop:22-fractional-kernel}
设整数 \(d\geq1\)、\(0<s<2\)，
\(c_{d,s}=(2\pi)^s/\int_{\mathbb R^d}(1-\cos(2\pi z_1))|z|^{-d-s}\,\mathrm dz\)，并以符号 \((2\pi|\xi|)^s\) 定义 \((-\Delta)^{s/2}\)。对复 Schwartz 函数 \(f\in\mathcal S(\mathbb R^d)\)，下面的第一式逐点绝对收敛，且
\[
 (-\Delta)^{s/2}f(x)=\frac{c_{d,s}}2
 \int_{\mathbb R^d}\frac{2f(x)-f(x+y)-f(x-y)}{|y|^{d+s}}\,\mathrm dy.
\]
以 \(\langle g,f\rangle=\int g\cdot\overline f\) 记复积分配对，还有
\[
 \langle(-\Delta)^{s/2}f,f\rangle
 =\frac{c_{d,s}}2
 \iint\frac{|f(x)-f(y)|^2}{|x-y|^{d+s}}\,\mathrm dx\,\mathrm dy.
\]
\end{proposition}
\begin{proof}
Taylor 积分公式使对称二阶差在原点为 \(O(|y|^2)\)，远场则由 \(f\) 的有界性控制。因为 \(0<s<2\)，这两端都可积，先得到逐点绝对收敛。还需使 Fourier 变换中的换序合法。对二阶差在 \(x\) 上取 \(L^1\) 范数，平移不变性与同一 Taylor 公式给出
\[
 \|2f-f(\cdot+y)-f(\cdot-y)\|_1
 \leq\min\left\{4\|f\|_1,\,
             |y|^2\sum_{j,k}\|\partial_j\partial_k f\|_1\right\}.
\]
该上界乘以 \(|y|^{-d-s}\) 后可积，Fubini 因而允许在全积分中作 Fourier 变换。所得乘子为
\[
 c_{d,s}\int_{\mathbb R^d}
       (1-\cos(2\pi y\cdot\xi))|y|^{-d-s}\,\mathrm dy.
\]
在 \(\xi\ne0\) 时先旋转，再作 \(z=|\xi|y\)，它等于
\(c_{d,s}|\xi|^s A_{d,s}=(2\pi|\xi|)^s\)；在 \(\xi=0\) 两边均为零。这证明第一式与 Fourier 定义相同。
两边都连续：Fourier 侧的符号乘 \(\widehat f\) 可积，实空间侧用有界二阶导数及有界函数值给出与 \(x\) 无关的可积控制。因此由 Fourier 变换得到的几乎处处等式实际逐点成立。

能量公式可直接用非负积分计算。Plancherel 给出
\[
 \|f(\cdot+y)-f\|_2^2
 =2\int_{\mathbb R^d}
        (1-\cos(2\pi y\cdot\xi))\cdot|\widehat f(\xi)|^2\,\mathrm d\xi.
\]
两边乘以 \(c_{d,s}|y|^{-d-s}/2\) 再积分。Tonelli 和刚算出的乘子把右边变为
\(\int(2\pi|\xi|)^s\cdot|\widehat f(\xi)|^2\,\mathrm d\xi\)；
左边令第二个位置变量等于 \(x+y\)，便是题中的双重积分。积分有限性也由
\(\|f(\cdot+y)-f\|_2^2\leq
\min\{4\|f\|_2^2,|y|^2\cdot\|\nabla f\|_2^2\}\)
直接保证。
\end{proof}
若 \(f\geq0\) 且紧支集，在支集外有
\((-\Delta)^{s/2}f(x)=-c_{d,s}\int f(y)|x-y|^{-d-s}\,\mathrm dy<0\)，除非 \(f=0\)。
因此这个导数一般不保持支集，也不保持非负性；它的二次型却非负。这三个性质并不矛盾：非负的是总能量，而不是每一点的算子输出。与整数 Laplace 算子的局部性相比，这个尾部正是分数阶模型的本质区别。

频率定义中的 Bessel 势也有一个不振荡的实空间核。取整数 \(d\ge1\)、\(s>0\)、\(x\in\mathbb R^d\setminus\{0\}\)，定义（原点值按可积函数的等价类解释）
\[
 G_s(x)=\frac1{\Gamma(s/2)}
 \int_0^\infty t^{s/2-1}e^{-t}
 (4\pi t)^{-d/2}e^{-|x|^2/(4t)}\,\mathrm dt.
\]
这是把热核作正权平均，权的总质量为一。它与齐次 Riesz 势的区别恰在额外因子 \(e^{-t}\)：大时间被截去，因而不会保留齐次核的长距离幂尾。
\Needspace{6\baselineskip}
\begin{proposition}[Bessel 位势的正核]\label{b2:prop:22-bessel-kernel}
设整数 \(d\geq1\)、\(s>0\)，定义
\[
 G_s(x)=\frac1{\Gamma(s/2)}
 \int_0^\infty t^{s/2-1}e^{-t}
 (4\pi t)^{-d/2}e^{-|x|^2/(4t)}\,\mathrm dt,
 \quad x\ne0.
\]
原点值按可积函数的等价类解释。令 \(J^{-s}\) 为符号 \((1+4\pi^2|\xi|^2)^{-s/2}\) 的乘子。有 \(G_s\geq0\)、\(\int G_s=1\)、\(\widehat G_s(\xi)=(1+4\pi^2|\xi|^2)^{-s/2}\)。故 \(J^{-s}f=G_s*f\)，并对全部 \(1\leq p\leq\infty\)、\(f\in L^p(\mathbb R^d;\mathbb C)\) 满足 \(\|J^{-s}f\|_p\leq\|f\|_p\)。
当 \(0<s<d\) 时，\(G_s(x)\asymp |x|^{s-d}\) 于 \(0<|x|\leq1\)；当 \(s=d\) 时量级为 \(1+\log(1/|x|)\)；当 \(s>d\) 时在原点附近有界。无穷远处三种情形均有指数衰减。
\end{proposition}
\begin{proof}
非负性与 Tonelli 给出总质量一。热核 Fourier 变换为 \(e^{-4\pi^2t|\xi|^2}\)，对 \(t\) 的 gamma 积分得到所述符号，Young 即给出压缩界。卷积定义对包括 \(p=\infty\) 的所有输入都有意义；其分布 Fourier 变换也满足该乘子公式。若 \(x\) 限制在远离零点的紧集，热核积分在 \(t\to0\) 的指数衰减控制任意次 \(x\) 导数，\(t\to\infty\) 则由 \(e^{-t}\) 控制，所以 \(G_s\) 在原点之外光滑。

记 \(r=|x|\)。在 \(0<r\leq1\) 时，积分的 \(t\geq1\) 部分有统一有限上界。对 \(0<t<1\)，用 \(e^{-1}\leq e^{-t}\leq1\)，再换元 \(u=r^2/(4t)\)，得到与
\[
 r^{s-d}\int_{r^2/4}^\infty
 u^{(d-s)/2-1}e^{-u}\,\mathrm du
\]
可比的表达。若 \(s<d\)，积分趋于一个正的 gamma 常数；若 \(s=d\)，零点附近是 \(\int_{r^2/4}^1\mathrm du/u\)；若 \(s>d\)，它的下限奇性恰好抵消外面的 \(r^{s-d}\)，结果有界。下界可以分别在固定 \(u\) 区间或 \(r^2/4<u<1\) 中取出，故这些不仅是上界。

当 \(r\geq1\)，使用 \(t+r^2/(4t)\geq r\)，有
\[
 e^{-t-r^2/(4t)}
 \leq e^{-r/2}e^{-t/2-r^2/(8t)}
 \leq e^{-r/2}e^{-t/2-1/(8t)}.
\]
最后一个因子乘上任意固定实数幂的 \(t\) 后，在零点及无穷远都可积。因此
\(G_s(x)\leq C_{d,s}e^{-|x|/2}\) 于 \(|x|\geq1\)，给出明确的远场上界。
\end{proof}
例如在 \(d=1,s=2\) 时，
\[
 G_2(x)=\frac12e^{-|x|}.
\]
事实上，右侧的 Fourier 变换为
\(\int_0^\infty e^{-x}\cos(2\pi x\xi)\,\mathrm dx
=(1+4\pi^2\xi^2)^{-1}\)，所以由 \(L^1\) Fourier 变换的单射性得到公式。它在原点连续，但一阶导数的右极限减左极限等于 \(-1\)，因而分布恒等式
\((1-\partial_x^2)G_2=\delta_0\) 成立。可积核与导数中的点质量同时出现，正好解释卷积解为什么满足二阶方程。

\ProseParagraphBreak

对任意 \(s,t>0\)，两个正核的卷积也属于 \(L^1\)，Fourier 变换相乘得到
\[
 G_s*G_t=G_{s+t}.
\]
因此连续施加两个 Bessel 位势把阶数相加。这与 \(J^{-s}J^{-t}=J^{-(s+t)}\) 一致，但这里还给出实空间的正卷积表示及端点压缩。

\ProseParagraphBreak

这个核解释了非齐次空间的低频项：常数频率附近的符号有界，所以 \(J^{-s}\) 保持同一 \(L^p\)；齐次符号 \(|\xi|^{-s}\) 在零频发散，改善可积性时必须改变目标指数或规定低频归一化。两者在高频上具有相同阶数，却不能在全空间上互相替代。

\ProseParagraphBreak

核估计还给出 Sobolev 嵌入的另一种表述。若 \(1<p<\infty\)、\(0<s<d/p\)，置 \(1/q=1/p-s/d\)。由 \(G_s(x)\leq C|x|^{s-d}\) 和 HLS，
\[
 \|f\|_q=\|G_s*(J^sf)\|_q\leq C\|J^sf\|_p.
\]
若 \(s>d/p\)，则 \(G_s\in L^{p'}\)：在 \(s<d\) 时，原点幂次条件正是
\((d-s)p'<d\)；\(s=d\) 时的对数在有限指数下可积，\(s>d\) 时局部有界，远场均由指数衰减处理。因此
\(\|f\|_\infty\leq\|G_s\|_{p'}\cdot\|J^sf\|_p\)，并得到连续代表。具体地，写 \(g=J^sf\)，则
\[
 |f(x+h)-f(x)|
 \leq\|G_s(\cdot+h)-G_s\|_{p'}\cdot\|g\|_p.
\]
有限指数下的核平移连续性使右端趋零，且与 \(x\) 无关。把 \(g\) 用有界紧支撑函数在 \(L^p\) 中逼近，卷积误差就在一致范数中趋零；每个近似卷积由核的远场指数界在无穷远趋零。因此该代表属于 \(C_0(\mathbb R^d)\)。

\ProseParagraphBreak

在临界 \(s=d/p\) 时，核不属于 \(L^{p'}\)，不能原样保留无穷端点。取
\[
 h(x)=|x|^{-s}\bigl(\log(e/|x|)\bigr)^{-\beta}
          \cdot\mathbf1_{\{\,0<|x|<1/2\,\}},
 \quad \frac1p<\beta\leq1.
\]
径向积分给出
\(\int|h|^p=C_d\int_0^{1/2}r^{-1}(\log(e/r))^{-\beta p}\,\mathrm dr<\infty\)。
于是 \(f=G_s*h\in H^{s,p}\)，并且 \(f\) 几乎处处有限。另一方面，正核的近原点下界给出
\[
 f(0)\geq c\int_0^{1/2}
         r^{-1}\bigl(\log(e/r)\bigr)^{-\beta}\,\mathrm dr=\infty.
\]
核在原点之外连续，所以对任意 \(x_n\to0\)，被积函数
\(G_s(x_n-y)\cdot h(y)\) 在几乎每个 \(y\) 趋于 \(G_s(-y)\cdot h(y)\)。
Fatou 因而给出 \(\liminf_n f(x_n)=\infty\)。这说明输出在原点的整个小邻域内可任意大，排除了通过修改一个零测集恢复本性有界性的可能。故临界 \(H^{d/p,p}\) 一般不嵌入 \(L^\infty\)。

\ProseParagraphBreak

最后，椭圆导数估计与这些嵌入处理不同问题。前者从方程右端恢复同指数的最高阶导数，后者把已经拥有的导数换成更高的可积指数或连续性。应用偏微分方程时，常见路线是先证明解处于 \(H^{s,p}\)，再按 \(s\) 与 \(d/p\) 的相对位置使用嵌入。若只看方程的阶数就声称解连续，会遗漏右端的可积指数及维数这两个决定性参数。

\ProseParagraphBreak

齐次导数与非齐次范数之间还有一个常用的精确比较。对 \(s>0\)、\(1<p<\infty\)，在 Schwartz 类上有
\[
 \|J^sf\|_p\asymp\|f\|_p+\|(-\Delta)^{s/2}f\|_p.
\]
正向用两个乘子：\(J^{-s}\) 以及
\((2\pi|\xi|)^s(1+4\pi^2|\xi|^2)^{-s/2}\)。
后者在原点不一定普通光滑，但定义域只要求 \(\xi\ne0\)，其尺度导数依然受控，故满足 Mihlin 条件。反向取低频截断 \(\chi\)，低频 \(J^s\chi\) 的核可积；高频符号
\[
 (1-\chi(\xi))\frac{(1+4\pi^2|\xi|^2)^{s/2}}{(2\pi|\xi|)^s}
\]
是零阶 Mihlin 符号，把它作用于齐次导数即可恢复高频 \(J^sf\)。两部分合并证明比较；Schwartz 类在 \(H^{s,p}\) 中稠密，因此按
\(\|f\|_p+\|(-\Delta)^{s/2}f\|_p\) 完备化得到同一个空间。这里的齐次导数是由这个图范数闭包定义的闭算子：若 \(f_n\to0\) 于 \(L^p\) 且齐次导数在 \(L^p\) 中收敛，范数比较使 \(J^sf_n\) 收敛；分布连续性迫使其极限为零，进而齐次导数极限也为零。它没有借此被定义到任意 \(\mathcal S'\) 上。

\ProseParagraphBreak

因此非齐次空间不是拒绝齐次导数，而是在其之外补上低阶 \(L^p\) 大小。高频上的两个定义等价，真正差别留在零频附近。使用能量方法时，可以在已说明的空间范围内自由改写这个图范数；使用全空间齐次逆、缩放不变估计或商掉多项式时，却必须保留齐次定义自己的归一化。一个范数比较不能自动解除所有齐次分布的低频歧义。

\subsection{椭圆型估计}
\label{b2:subsec:220303}

从方程右端恢复二阶导数时，要除以二阶主符号。若它沿某个方向为零，除法可能失去控制；正定性则给出与方向无关的分母下界。先处理常系数全空间问题，随后说明低频项、边界和系数变化分别增加什么困难。本小节记
\[
 \|D^2u\|_p\coloneqq\sum_{j,k=1}^d\|\partial_j\partial_ku\|_p.
\]

\Needspace{6\baselineskip}
\begin{proposition}[常系数椭圆导数估计]\label{b2:prop:22-elliptic-estimate}
设整数 \(d\geq1\)、\(1<p<\infty\)，\(A=(a_{jk})\) 是实对称常矩阵，并且存在 \(0<\nu\leq\Lambda<\infty\)，使
\[
 \nu|\xi|^2\leq\xi^{\mathsf T}A\xi\leq\Lambda|\xi|^2
 \quad(\xi\in\mathbb R^d).
\]
令 \(Lu\coloneqq-\sum_{j,k}a_{jk}\partial_j\partial_ku\)。则每个
\(u\in W^{2,p}(\mathbb R^d;\mathbb C)\) 满足
\[
 \|D^2u\|_p\leq C_{d,p,\nu,\Lambda}\|Lu\|_p.
\]
所有导数及方程均按分布理解；对光滑函数，它们与通常导数相同。
\end{proposition}
\begin{proof}
先设 \(u\in C_c^\infty\)。记 \(q(\xi)=\xi^{\mathsf T}A\xi\)，则
\[
 \widehat{Lu}(\xi)=4\pi^2q(\xi)\cdot\widehat u(\xi),\quad
 \widehat{\partial_j\partial_ku}(\xi)
 =-\frac{\xi_j\xi_k}{q(\xi)}\widehat{Lu}(\xi)
 \quad(\xi\ne0).
\]
商 \(m_{jk}=-\xi_j\xi_k/q\) 零次齐次。在单位球面上，\(q\geq\nu\)，而 \(q\) 的至多二阶导数由 \(\Lambda\) 控制，更高阶导数为零。商的反复求导只产生这些导数的有限乘积及 \(q\) 的负幂，故单位球面上的至多 \(d+2\) 阶导数由 \(d,\nu,\Lambda\) 一致控制。齐次性再给出
\[
 |\partial^\alpha m_{jk}(\xi)|
 \leq C_{d,\alpha,\nu,\Lambda}|\xi|^{-|\alpha|}.
\]
由\cref{b2:thm:22-mihlin}得到每个二阶导数的界，有限相加即可。

对一般 \(W^{2,p}\) 函数，先以光滑单位积分核磨光，其至多二阶弱导数均在 \(L^p\) 中收敛。再乘以在 \(B_R\) 为一、在 \(B_{2R}\) 外为零的缩放截断。乘积求导产生的附加项由
\(CR^{-1}\|\nabla u\|_p\)、\(CR^{-2}\|u\|_p\) 控制，原有函数及导数在 \(|x|>R\) 上的尾范数趋零。因此可以取得 \(u_n\in C_c^\infty\)，使 \(u_n\to u\) 于 \(W^{2,p}\)。常系数保证 \(Lu_n\to Lu\) 于 \(L^p\)，将已证估计传给极限即得结论。
\end{proof}
这个命题是先验估计：它从已经存在的 \(u\) 及 \(Lu\) 比较范数，没有构造对每个右端 \(f\) 的逆解。即使 \(A=I\)，低频仍阻止同指数的无质量逆。取非零 \(\psi\in C_c^\infty\)、\(u_R(x)=\psi(x/R)\)，则
\[
 \|u_R\|_p=R^{d/p}\|\psi\|_p,\quad
 \|Lu_R\|_p=R^{d/p-2}\|L\psi\|_p.
\]
让 \(R\to\infty\)，便排除独立于 \(u\) 的估计 \(\|u\|_p\leq C\|Lu\|_p\)。二阶导数在相同缩放下也带有 \(R^{-2}\)，所以刚证明的最高阶估计没有矛盾。

\ProseParagraphBreak

添加正的零阶项以后，频率原点的除法也成为合法操作。这个变化既给出解的存在唯一性，也产生一个尺度一致的参数族。

\Needspace{6\baselineskip}
\begin{proposition}[带正参数的椭圆预解算子]\label{b2:prop:22-elliptic-resolvent}
设整数 \(d\geq1\)、\(1<p<\infty\)，实对称常矩阵 \(A\) 满足
\(\nu|\xi|^2\leq\xi^{\mathsf T}A\xi\leq\Lambda|\xi|^2\)，其中
\(0<\nu\leq\Lambda<\infty\)。令
\(L=-\sum_{j,k}a_{jk}\partial_j\partial_k\)。对每个 \(\lambda>0\) 和
\(f\in L^p(\mathbb R^d;\mathbb C)\)，分布方程
\[
 (\lambda+L)u=f
\]
有唯一的温和分布解。该解属于 \(W^{2,p}(\mathbb R^d)\)，记为
\(u=(\lambda+L)^{-1}f\)，并满足
\[
 \lambda\|u\|_p+\sqrt\lambda\|\nabla u\|_p+\|D^2u\|_p
 \leq C_{d,p,\nu,\Lambda}\|f\|_p.
\]
常数与 \(\lambda,f\) 无关。特别地，\((1+L)^{-1}:L^p\to W^{2,p}\) 是有界解算子。
\end{proposition}
\begin{proof}
函数 \(b_\lambda(\xi)=(\lambda+4\pi^2\xi^{\mathsf T}A\xi)^{-1}\) 在全空间光滑，且对固定 \(\lambda>0\)，它及其各阶导数至多多项式增长。因此可在 \(\mathcal S'\) 中定义
\(\widehat u=b_\lambda\cdot\widehat f\)，并立即得到分布方程。若齐次方程的温和分布解为 \(v\)，把
\((\lambda+4\pi^2\xi^{\mathsf T}A\xi)\widehat v=0\)
乘以同一个光滑逆，即得 \(\widehat v=0\)。这证明的是整个温和分布类中的唯一性。

为证明正则性与统一界，考察三个加权符号
\[
 \frac{\lambda}{\lambda+4\pi^2\xi^{\mathsf T}A\xi},\quad
 \frac{\sqrt\lambda\,2\pi\mathrm i\xi_j}
      {\lambda+4\pi^2\xi^{\mathsf T}A\xi},\quad
 \frac{-4\pi^2\xi_j\xi_k}
      {\lambda+4\pi^2\xi^{\mathsf T}A\xi}.
\]
令 \(\xi=\sqrt\lambda\,\zeta\)，它们分别变为
\[
 \frac1{1+4\pi^2\zeta^{\mathsf T}A\zeta},\quad
 \frac{2\pi\mathrm i\zeta_j}{1+4\pi^2\zeta^{\mathsf T}A\zeta},\quad
 \frac{-4\pi^2\zeta_j\zeta_k}{1+4\pi^2\zeta^{\mathsf T}A\zeta}.
\]
在 \(|\zeta|\leq1\) 时，分母至少为一，有限阶导数有界。在
\(|\zeta|\geq1\) 时，分母至少为 \(4\pi^2\nu|\zeta|^2\)；反复求导分别给出阶数 \(-2,-1,0\) 的符号界。因此三者都有只依赖
\(d,\nu,\Lambda\) 的 Mihlin 常数。尺度导数中的
\(|\xi|^{|\alpha|}\) 正好抵消缩放微分产生的
\(\lambda^{-|\alpha|/2}\)，所以原来的参数族有相同常数。

对 Schwartz 右端应用 Mihlin 定理，就得到题中估计。对一般
\(f\in L^p\)，取 Schwartz 逼近。该估计使对应 \(u_n\) 及一、二阶导数分别在 \(L^p\) 中收敛；分布求导与极限相容，所以极限属于 \(W^{2,p}\) 并满足方程。前面已证的分布唯一性将它与 Fourier 定义的 \(u\) 识别，统一界也传给极限。
\end{proof}
取 \(A=I,\lambda=1\)，解就是前面的 \(J^{-2}f=G_2*f\)。在一维还可直接写为
\[
 u(x)=\frac12\int_{\mathbb R}e^{-|x-y|}\cdot f(y)\,\mathrm dy.
\]
该公式把正核、分布中的导数跳跃与 \(W^{2,p}\) 解连接起来。对于一般对角矩阵
\(A=\operatorname{diag}(a_1,\ldots,a_d)\)，换元
\(x_j=\sqrt{a_j}\,y_j\) 将 \(L\) 化为 \(-\Delta_y\)，而
\(\partial_{x_j}=a_j^{-1/2}\partial_{y_j}\)。在一维，
\(\|u''\|_p=a_1^{-1}\|Lu\|_p\) 恰为等式，表明椭圆常数趋零时最高阶估计确会退化。一般正定矩阵先作固定正交对角化即可理解同一现象。

\ProseParagraphBreak

当 \(\lambda\downarrow0\)，最高阶导数界仍然统一，但零阶估计只给出
\(\|u\|_p\leq C\lambda^{-1}\|f\|_p\)。这没有产生同指数的 \(L^{-1}\)。温和分布中的无质量齐次方程还允许满足 \(LP=0\) 的多项式；其中仿射函数的二阶导数为零，而在 \(d\geq2\) 时，更高次的这类多项式可以具有非零二阶导数。这些非零多项式不属于本节有限指数的 \(W^{2,p}(\mathbb R^d)\)，不能代入先验命题。自由度来自所选解空间及低频归一化，不能统称为二阶导数都看不见的项。

\ProseParagraphBreak

全空间与有边界区域也要区别。在 \(\mathbb R^2\) 的单位球上，
\(u(x)=x_1x_2\) 满足 \(\Delta u=0\)，却有
\(\partial_1\partial_2u=1\)。所以不加边界条件，不能把全空间结论改写为
\(\|D^2u\|_{L^p(B_1)}\leq C\|\Delta u\|_{L^p(B_1)}\)。
若只求内部估计，截断后的低阶项会记录边界之外的信息。

\ProseParagraphBreak

变系数时，空间乘法不与 Fourier 变换交换，不能把 \(a_{jk}(x)\) 当成常数塞进上述商。下面的冻结步骤说明，在一个小球内接近固定矩阵时，哪一部分误差可以吸收。这里选择非散度形式，因而不要求对系数求导。

\Needspace{6\baselineskip}
\begin{proposition}[冻结系数的先验估计]\label{b2:prop:22-freezing-estimate}
设整数 \(d\geq1\)、\(1<p<\infty\)、\(r>0\)，以
\(\|D^2v\|_p=\sum_{j,k=1}^d\|\partial_j\partial_kv\|_p\)
记二阶导数和范数。令 \(A_0\) 为实对称正定常矩阵，并取常数 \(C_0\) 使
\[
 \|D^2v\|_p\leq C_0\|L_{A_0}v\|_p
 \quad(v\in W^{2,p}(\mathbb R^d)),\quad
 L_{A_0}v=-\sum_{j,k}a_{jk}^0\partial_j\partial_kv.
\]
设实矩阵场 \(A=(a_{jk})\) 在 \(B_r\subset\mathbb R^d\) 上有界可测，满足
\[
 \max_{j,k}\|a_{jk}-a_{jk}^0\|_{L^\infty(B_r)}
 \leq\varepsilon,\quad C_0\varepsilon\leq\frac12.
\]
对 \(u\in W^{2,p}(B_r)\)，以可测函数
\(L_Au=-\sum_{j,k}a_{jk}\partial_j\partial_ku\) 定义非散度算子。
若 \(u\) 的支撑紧包含于 \(B_r\)，则
\[
 \|D^2u\|_{L^p(B_r)}\leq2C_0\|L_Au\|_{L^p(B_r)}.
\]
一般 \(u\in W^{2,p}(B_r)\) 则满足
\[
 \begin{aligned}
 \|D^2u\|_{L^p(B_{r/2})}
 \leq C\bigl(&\|L_Au\|_{L^p(B_r)}
       +r^{-1}\|\nabla u\|_{L^p(B_r)}\\
       &+r^{-2}\|u\|_{L^p(B_r)}\bigr),
 \end{aligned}
\]
其中 \(C\) 可依赖 \(d,p,C_0\) 及各系数的 \(L^\infty\) 上界，但与 \(u,r\) 无关。
\end{proposition}
\begin{proof}
先设 \(u\) 紧支于球内。它在球外延零后属于全空间 \(W^{2,p}\)：在边界的一个邻域内原函数已经为零，所以分布求导没有边界项。由系数差的上界及本小节选定的二阶导数和范数，
\[
 \|(L_{A_0}-L_A)u\|_p
 \leq\varepsilon\sum_{j,k}\|\partial_j\partial_ku\|_p
 =\varepsilon\|D^2u\|_p.
\]
常系数命题于是给出
\[
 \|D^2u\|_p
 \leq C_0\|L_Au\|_p+C_0\varepsilon\|D^2u\|_p.
\]
移项即得第一式。

对一般 \(u\)，取 \(\eta\in C_c^\infty(B_r)\)，在 \(B_{r/2}\) 上为一，并使
\(\|\nabla\eta\|_\infty\leq C_dr^{-1}\)、
\(\|D^2\eta\|_\infty\leq C_dr^{-2}\)。弱导数的乘积公式给出
\[
 \begin{aligned}
 L_A(\eta\cdot u)
 ={}&\eta\cdot L_Au
 -\sum_{j,k}a_{jk}
       \bigl((\partial_j\eta)\cdot\partial_ku
              +(\partial_k\eta)\cdot\partial_ju\bigr)\\
 &-\sum_{j,k}a_{jk}(\partial_j\partial_k\eta)\cdot u.
 \end{aligned}
\]
对 \(\eta\cdot u\) 使用第一式，再以系数上界及截断导数界控制右端。在
\(B_{r/2}\) 内，\(D^2(\eta\cdot u)=D^2u\)，因而得到内部估计。
\end{proof}
连续系数在足够小的球上可以满足该小振荡条件；结论的输入仍是已经属于 \(W^{2,p}\) 的函数。若起始解只有 \(H^1\)，尚不能先把
\(a_{jk}\partial_j\partial_ku\) 当作 \(L^p\) 函数，再以这个估计证明二阶导数存在。存在性或由弱解提升正则性还需要合法的逼近、解的构造及收敛论证。

\ProseParagraphBreak

散度算子是另一个对象。只有系数有足够的弱导数时，才可把
\(-\sum_{j,k}\partial_j(a_{jk}\partial_ku)\) 展开成非散度部分及
\(-\sum_{j,k}(\partial_ja_{jk})\partial_ku\)。任意有界可测系数不允许这样把额外项当作普通函数移到右侧。\chapref{b2:ch:16}和\appref{b2:app:H}分别处理相应椭圆正则问题及低正则系数；乘子估计则为全空间常系数方程及小振荡扰动提供先验控制。

\subsection{球面乘子 预告}
\label{b2:subsec:220304}

径向符号只依赖 \(|\xi|\)，仍可能在球面上发生跳跃。
光滑环带截断 \(\chi(|\xi|)\) 的核属于 \(L^1\)；球特征函数 \(\mathbf1_{\{|\xi|\le1\}}\) 则不满足任何跨边界的 Mihlin 光滑条件。
把跳跃平滑成宽度 \(\epsilon\) 的过渡层时，第 \(N\) 阶导数通常增长如 \(\epsilon^{-N}\)。对每个固定 \(\epsilon\) 应用定理可得有界性，却得不到 \(\epsilon\downarrow0\) 的统一界。

\ProseParagraphBreak
Bochner–Riesz 平均以 \((1-|\xi|^2)_+^\delta\) 逐渐消减边界跳跃，其所需光滑阶数与 \(p\)、维数有关。Mihlin 定理给出方便的充分条件，但常不是锐利范围。\chapref{b2:ch:24}用径向核的振荡和曲率研究这一问题；\chapref{b2:ch:23}的限制估计控制频率薄壳上的集中，则说明为什么乘子问题会与曲面延拓相遇。

\ProseParagraphBreak
因此“径向”“有界”“紧支撑”三项都不能单独替代乘子条件。对核使用 Young 不等式、对符号使用尺度导数、对薄壳使用曲面几何，是三条不同证明路线。应用时应明确正在使用哪一条，以及它在端点或退化极限中会失去哪一个条件。

平滑乘子还有一个时间演化模型，可与频率硬截断比较。对整数 \(d\geq1\)、\(t>0\)，热算子 \(e^{t\Delta}\) 的符号是 \(e^{-4\pi^2t|\xi|^2}\)，核为
\[
 H_t(x)=(4\pi t)^{-d/2}e^{-|x|^2/(4t)}.
\]
其空间核非负，积分为一，因而在全部 \(L^p\) 上压缩。写
\(H_t(x)=t^{-d/2}H_1(x/\sqrt t)\)，对多重指标 \(\alpha\) 求导得到
\[
 \partial^\alpha H_t(x)
 =t^{-(d+|\alpha|)/2}
       (\partial^\alpha H_1)(x/\sqrt t),\quad
 \|\partial^\alpha H_t\|_1
 =t^{-|\alpha|/2}\|\partial^\alpha H_1\|_1.
\]
固定时间核的每个导数都是多项式乘 Gaussian，故右端常数有限。Young 不等式于是对每个整数 \(k\geq0\) 给出
\[
 \|D^ke^{t\Delta}f\|_p\leq C_k t^{-k/2}\|f\|_p,
 \quad 1\leq p\leq\infty,\quad t>0.
\]
这里 \(\|D^k u\|_p\) 可取全部 \(k\) 阶导数范数的有限和。这个证明连两个端点也包括。每个正时间的任意阶导数还可由核的平移连续性得到连续代表，因而产生光滑函数。单纯的开区间 Mihlin 定理已经足以处理 \(1<p<\infty\)，但核方法在这个模型中提供了额外端点。

\ProseParagraphBreak

任意实 \(s\geq0\) 时同样有
\[
 \|e^{t\Delta}f\|_{H^{s,p}}\leq C_{s,p}(1+t^{-s/2})\|f\|_p,
 \quad1<p<\infty.
\]
在 \(0<t\leq1\) 中，把 \(t^{s/2}(1+4\pi^2|\xi|^2)^{s/2}e^{-4\pi^2t|\xi|^2}\) 以 \(\zeta=\sqrt t\,\xi\) 重写，成为
\((t+4\pi^2|\zeta|^2)^{s/2}e^{-4\pi^2|\zeta|^2}\)。
为看清常数为何统一，先只对幂函数求 \(\alpha\) 阶导数。每一项都由
\[
 C_{\alpha,\beta,\ell}\,
 \zeta^\beta(t+4\pi^2|\zeta|^2)^{s/2-\ell},
 \quad |\beta|=2\ell-|\alpha|
\]
组成。乘上 \(|\zeta|^{|\alpha|}\)，再用
\(|\zeta|^{2\ell}\leq C_\ell(t+4\pi^2|\zeta|^2)^\ell\)，就控制为
\(C(t+4\pi^2|\zeta|^2)^{s/2}\)。乘积求导中落到 Gaussian 的部分只增加固定多项式，在无穷处由 Gaussian 衰减吸收；在原点处没有剩下负幂。因此 \(0<t\leq1\) 的 Mihlin 常数统一。若 \(t\geq1\)，则把一次固定时间的平滑与剩余时间的压缩复合；\(t=1\) 时后者是恒等算子。这说明分数阶平滑不必先把 \(s\) 取整。

\ProseParagraphBreak

当 \(t\downarrow0\)，正阶范数常数发散，与一般 \(L^p\) 数据可能没有任何正阶导数完全相容。对每个固定正时间得到光滑函数，不意味着这些函数在更强空间中一致趋向原数据。对 \(1\leq p<\infty\)，热核的近似恒等性保证在原来的 \(L^p\) 中趋于初值，提升正则性则付出时间尺度。一般 \(L^\infty\) 输入没有这种范数强连续性：取半空间示性函数 \(f=\mathbf1_{\{x_1>0\}}\)，热卷积在边界附近连续趋于 \(1/2\)，原函数在两侧分别为零和一，因此每个 \(t>0\) 都有 \(\|e^{t\Delta}f-f\|_\infty\geq1/2\)。

\ProseParagraphBreak

把时间倒过来，固定 \(t>0\) 后符号成为 \(e^{4\pi^2t|\xi|^2}\)，甚至失去 \(L^2\) 乘子所需的本性有界性。取非零 \(g\in\mathcal S\)，使
\(\widehat g\in C_c^\infty(B(0,1))\)，并令
\(f_N(x)=e^{2\pi\mathrm iNx_1}\cdot g(x)\)、\(N>1\)。
其 Fourier 支撑落在 \(B(Ne_1,1)\)，所以对这些频率紧支撑输入，逆热像仍可定义，且 Plancherel 给出
\[
 \|f_N\|_2=\|g\|_2,\quad
 \|e^{-t\Delta}f_N\|_2
 \geq e^{4\pi^2t(N-1)^2}\|g\|_2.
\]
再取 \(h_N=e^{-2\pi^2t(N-1)^2}f_N\)，便有 \(\|h_N\|_2\to0\)，而
\(\|e^{-t\Delta}h_N\|_2\to\infty\)。正向平滑因此不能通过形式上的指数取反得到稳定逆过程。

\ProseParagraphBreak

与球面乘子相比，热符号没有一条频率边界上的跳跃，所有尺度的导数也具有统一控制；与齐次分数积分相比，它在零频不发散。良好的端点压缩、正时间平滑和低频稳定是这些特殊性质共同产生的，不应当归给任意一个写成 \(m(\xi)\) 的符号。

\ProseParagraphBreak

频率过渡层还可以用一个明确的算子范数检验。取固定 \(\chi\in C_c^\infty\)，令 \(m_R(\xi)=\chi(\xi/R)\)。其核为 \(R^d\check\chi(Rx)\)，所以全部 \(R>0\) 的 \(L^1\) 核范数都等于 \(\|\check\chi\|_1\)，在每个 \(L^p\) 上有同一个界。若 \(\chi(0)=1\)，在 \(R\to\infty\) 时，有限指数下这族算子强趋向恒等：先对 Schwartz 输入用核的近似恒等性，再用一致范数和稠密性推广。将频率支撑推向无穷本身并不造成常数发散；关键是过渡轮廓经过尺度归一化后仍固定。

\ProseParagraphBreak

若改用
\[
 m_\epsilon(\xi)=\theta\bigl((|\xi|-1)/\epsilon\bigr),
 \quad0<\epsilon<\frac12,
\]
其中 \(\theta\in C^\infty(\mathbb R)\)、\(0\leq\theta\leq1\)，在
\((-\infty,-1]\) 为一、在 \([1,\infty)\) 为零，则过渡层紧贴单位球面。符号在原点附近恒为一，故原点处也光滑；但归一化整个频率空间不能同时固定球的半径和层的厚度。在某个固定的过渡参数处 \(\theta'\neq0\)，径向一阶导数就达到
\(c\epsilon^{-1}\)，而 \(|\xi|\) 仍约为一，尺度导数条件因此真实增长。前一个模型只有整体半径，后一个还带有厚度与半径之比；两者的参数一致性不同。

\ProseParagraphBreak

在 \(L^2\) 中只需符号的本性上确界，故这族越来越薄的过渡层始终满足
\(\|T_{m_\epsilon}\|_{2\to2}\leq1\)。对 Schwartz 输入，支配收敛还给出
\[
 T_{m_\epsilon}f\longrightarrow
 T_{\mathbf1_{\{|\xi|\leq1\}}}f\quad\text{于 }L^2.
\]
若某个 \(1<p<\infty\) 上另外已有与 \(\epsilon\) 无关的算子界，就可用\cref{b2:prop:22-multiplier-calculus}的极限结论，把该界传给球乘子。仅有每个固定 \(\epsilon\) 的 Mihlin 界不能完成这一步；导数常数增长本身也不证明极限算子一定无界。后两章将利用频率边界的几何进一步分析实际有界范围。


\begin{chaptersummary}
乘子问题先有一个明确的 \(L^2\) 算子，再问它能否保持其他空间。复合、对偶、坐标变化与受控极限提供了基本运算，但每一步都要保留指数、定义域和常数。有限指数下的稠密性与无穷指数下的伴随定义承担不同职责；逐个输入的强收敛，也不能替代整个单位球上的范数收敛。

\ProseParagraphBreak

Mihlin 与 Hörmander 条件都是可核验的充分条件。它们通过固定环带上的统一估计控制各尺度之间的叠加，再由奇异积分理论得到 \(1<p<\infty\) 的结论。端点若要成立，常常还须使用正核、有限测度或可积核等额外结构。某个导数判据失效，只说明该判据不能直接使用，不能据此断言乘子本身无界。

\ProseParagraphBreak

在应用中，Riesz 变换组织齐次导数，Bessel 位势补上低频大小，常系数椭圆符号恢复同指数的高阶导数。热核的正性与可积性则提供额外的端点和平滑性质。符号的阶数、原点行为、边界几何和参数尺度应分别核对；将这些信息合在一个没有说明的“光滑符号”条件中，会掩盖真正使用的假设。
\end{chaptersummary}

\BookExercises

% 自拟；基础核验与反例：有限原子卷积在不同指数上的准确范数，区分总变差、符号大小和可积核。
\Needspace{6\baselineskip}
\begin{exercise}\label{b2:ex:22-new-1}
在 \(\mathbb R\) 上令 \(\tau_af(x)=f(x-a)\)，并定义
\[
 Tf=f+\tau_1f-\tau_2f.
\]
求其 Fourier 符号，证明
\[
 \|T\|_{L^1\to L^1}=3,
 \quad \|T\|_{L^2\to L^2}=\sqrt5,
 \quad \|T\|_{L^\infty\to L^\infty}=3.
\]
说明为什么卷积核的总变差只给出一般上界，以及该算子为何不能由一个普通 \(L^1\) 卷积核表示。
\end{exercise}
\begin{solution}
平移公式给出
\[
 m(\xi)=1+e^{-2\pi\mathrm i\xi}-e^{-4\pi\mathrm i\xi}.
\]
三个平移都是等距映射，三角不等式给出所有 \(1\leq p\leq\infty\) 上的范数上界三。为求 \(L^1\) 下界，取非负 \(g\in C_c^\infty\bigl((-1/4,1/4)\bigr)\)，使 \(\|g\|_1=1\)。三个平移后的支撑互不相交，因此 \(\|Tg\|_1=3\)。

令 \(I=(-1/4,1/4)\)，取
\(f=\mathbf1_I+\mathbf1_{I-1}-\mathbf1_{I-2}\)。三个区间互不相交，\(\|f\|_\infty=1\)，而对 \(x\in I\) 有 \(Tf(x)=3\)。这在一个正测度集合上达到三，证明了本性上确界意义下的 \(L^\infty\) 下界。

令 \(z=e^{-2\pi\mathrm i\xi}\)，直接展开有
\[
 |1+z-z^2|^2=3-z^2-z^{-2}=3-2\cos(4\pi\xi).
\]
该连续周期函数的本性上确界为五。由\cref{b2:prop:22-l2}，\(L^2\) 范数恰为 \(\sqrt5\)。分布核为 \(\delta_0+\delta_1-\delta_2\)，其总变差为三，却大于这个平方空间范数。

若存在 \(k\in L^1\) 给出同一算子，则 \(\widehat k=m\) 几乎处处。两边连续，故处处相等；Riemann--Lebesgue 引理要求 \(\widehat k\) 在无穷远趋零，但 \(m(n)=1\) 对所有整数 \(n\) 成立，矛盾。有限原子核可以定义有界卷积算子，并不需要成为普通可积函数。
\end{solution}

% 自拟；中等：虚幂乘子群的参数导数界、强连续性与不范数连续，超出单一Mihlin验证。
\Needspace{6\baselineskip}
\begin{exercise}\label{b2:ex:22-new-2}
设整数 \(d\geq1\)、\(1<p<\infty\)，对 \(t\in\mathbb R\) 和 \(\xi\ne0\) 令
\[
 m_t(\xi)=|\xi|^{\mathrm it}
       =e^{\mathrm it\log|\xi|},
 \quad U_t=T_{m_t}.
\]
证明它们在 \(L^p(\mathbb R^d)\) 上形成一个由有界算子组成的一参数群，并有
\[
 \|U_t\|_{p\to p}\leq C_{d,p}(1+|t|)^{d+2}.
\]
再证明 \(t\mapsto U_tf\) 对每个 \(f\in L^p\) 都范数连续，但当 \(t\ne0\) 时
\[
 \|U_t-I\|_{p\to p}\geq2.
\]
说明这两个结论为何不矛盾，并求 \(p=2\) 时最后一个范数的准确值。原点处的符号值任取，不影响问题。
\end{exercise}
\begin{solution}
求一次导数得到
\(\partial_jm_t=\mathrm it\xi_j|\xi|^{-2}\cdot m_t\)。继续求导，每个 \(r\) 阶导数都可写成 \(|\xi|^{\mathrm it-r}\) 乘以关于 \(t\) 次数不超过 \(r\) 的多项式，其系数是单位球面上的光滑函数。因此
\[
 |\partial^\alpha m_t(\xi)|
 \leq C_\alpha(1+|t|)^{|\alpha|}|\xi|^{-|\alpha|}.
\]
取到 \(d+2\) 阶并应用\cref{b2:thm:22-mihlin}，得到题中的界。点态恒等式 \(m_{s+t}=m_sm_t\)、\(m_0=1\) 与乘子复合规则给出
\(U_{s+t}=U_sU_t\)、\(U_0=I\)、\(U_t^{-1}=U_{-t}\)。这里每个参数对应一个有界算子，并未断言全体参数的范数具有同一个上界。

先证明在零点强连续。对 \(f\in\mathcal S\)，有 \(m_t\to1\) 几乎处处，且 \(|m_t-1|\leq2\)，所以 Plancherel 与支配收敛给出
\(\|(U_t-I)f\|_2\to0\)。若 \(p\ne2\)，选 \(q\in(1,\infty)\)，使 \(p\) 严格介于 \(2\) 和 \(q\) 之间。对 \(|t|\leq1\)，已证乘子界使 \(\|(U_t-I)f\|_q\) 一致有界。空间范数的 Hölder 插值便使 \(\|(U_t-I)f\|_p\to0\)。再用 Schwartz 类在 \(L^p\) 中的稠密性及 \(|t|\leq1\) 的一致算子界，把结论推广到所有 \(f\in L^p\)。群性质将零点连续性移到任意参数点。

当 \(t\ne0\) 时，\(t\log r\) 随 \(r>0\) 走遍实轴。因此 \(|e^{\mathrm it\log r}-1|\) 的上确界为二，而且接近二的值出现在正测度环带上，故
\(\|m_t-1\|_\infty=2\)。由\cref{b2:prop:22-symbol-bound}，
\[
 \|U_t-I\|_{p\to p}\geq\|m_t-1\|_\infty=2.
\]
当 \(p=2\) 时等号成立。强连续固定一个输入再令参数趋零；算子范数则对整个单位球取上确界，所测试的频率范围可以随参数改变，两种收敛要求不同。
\end{solution}

% 自拟；基础到中等：带非零下界的Mihlin符号求逆，展开倒数导数而不假定算子范数小扰动。
\Needspace{6\baselineskip}
\begin{exercise}\label{b2:ex:22-new-3}
设整数 \(d\geq1\)，\(N=d+2\)，\(m\in C^N(\mathbb R^d\setminus\{0\};\mathbb C)\)，并满足
\[
 |m(\xi)|\geq\delta>0,
 \quad
 \sup_{\xi\ne0}|\xi|^{|\alpha|}\cdot
           |\partial^\alpha m(\xi)|\leq A
 \quad(|\alpha|\leq N).
\]
证明对每个 \(1<p<\infty\)，\(T_m\) 是 \(L^p\) 上的有界可逆算子，逆为 \(T_{1/m}\)，并给出仅依赖 \(d,p,A,\delta\) 的逆范数上界。可证明倒数的尺度导数统一不超过
\[
 C_N\delta^{-1}\sum_{j=0}^N(A/\delta)^j.
\]
将结果用于 \(m(\xi)=2+e^{\mathrm it\log|\xi|}\)、\(t\in\mathbb R\)，说明为何无需先证明 \(\|T_m-2I\|_{p\to p}<2\)。
\end{exercise}
\begin{solution}
记 \(n=1/m\)。零阶界为 \(|n|\leq\delta^{-1}\)。若 \(|\alpha|=r\geq1\)，反复微分倒数，得到有限项之和；每一项具有形式
\[
 c\,m^{-j-1}\prod_{\ell=1}^j\partial^{\beta_\ell}m,
 \quad
 1\leq j\leq r,
 \quad |\beta_\ell|\geq1,
 \quad\sum_{\ell=1}^j\beta_\ell=\alpha.
\]
该形式可由一次求导归纳：导数落到 \(m^{-j-1}\) 上时增加一个一阶导数因子，落到已有因子上时增加其多重指标。乘以 \(|\xi|^r\) 后，每一项至多为 \(C\delta^{-j-1}A^j\)。对有限项求和，并包括零阶界，即得到所给统一控制。

Mihlin 定理于是给出
\[
 \|T_{1/m}\|_{p\to p}
 \leq C_{d,p}\delta^{-1}\sum_{j=0}^{d+2}(A/\delta)^j.
\]
符号乘积 \(m\cdot(1/m)=1\) 几乎处处成立，乘子复合规则便说明
\(T_mT_{1/m}=T_{1/m}T_m=I\)。这同时给出每个右端的存在性和解的唯一性，而不只是一条先验不等式。

在给定例中，三角不等式给出 \(|m|\geq1\)。虚幂符号的导数计算又给出有限的 Mihlin 常数，故每个实 \(t\) 都有上述逆。算子级的几何级数需要一个小范数条件，本题直接控制倒数符号的导数，因而不需要验证该条件。结论使用了完整的 Mihlin 假设；这里只凭一个一般乘子符号远离零，并没有证明其倒数必属于同一个乘子空间。
\end{solution}

% 自拟；中等：张量乘子的准确范数乘积，推广一维投影并构造端点失败的Schwartz测试输入。
\Needspace{6\baselineskip}
\begin{exercise}\label{b2:ex:22-new-4}
设整数 \(d\geq1\)、\(1\leq p<\infty\)，且 \(m_j\in M_p(\mathbb R)\)、\(1\leq j\leq d\)。证明
\[
 m(\xi)=\prod_{j=1}^dm_j(\xi_j)
 \quad\Longrightarrow\quad
 \|m\|_{M_p(\mathbb R^d)}
       =\prod_{j=1}^d\|m_j\|_{M_p(\mathbb R)}.
\]
对上界按坐标使用 Fubini，对下界使用乘积输入。由此证明任意有限边长的矩形
\(Q=\prod_{j=1}^d[a_j,b_j]\)、\(a_j<b_j\)，其特征函数都是 \(1<p<\infty\) 上的乘子，且范数与各边长、中心无关。最后证明它不是 \(L^1\) 乘子，也没有规范的\WeakStar{}连续 \(L^\infty\) 乘子延拓。
\end{exercise}
\begin{solution}
先在有限和的乘积输入上，按各坐标先后应用一维算子。每一步由 Fubini 和一维 \(L^p\) 界控制，得到乘积上界；这类输入在有限指数 \(L^p(\mathbb R^d)\) 中稠密，所以算子唯一延拓。在同时属于 \(L^2\) 的乘积输入上，逐坐标 Fourier 变换说明该算子的符号正是 \(m\)，故所得延拓与乘子定义一致。

若各个范数均非零，对任意 \(\varepsilon>0\)，可以取单位输入 \(f_j\in\mathcal S(\mathbb R)\)，使
\(\|T_{m_j}f_j\|_p\geq\|m_j\|_{M_p}-\varepsilon\)。这是由 Schwartz 稠密性和有界性得到的。令 \(f(x)=\prod_jf_j(x_j)\)，则 \(\|f\|_p=1\)，而输出的范数恰为各个一维输出范数的乘积。令 \(\varepsilon\downarrow0\) 得下界。若某个范数为零，则对应符号为零几乎处处，等式仍成立。

一维区间投影已在正文由 Hilbert 变换建立。设 \(C_p=\|\mathbf1_{[0,1]}\|_{M_p(\mathbb R)}\)，频率平移和缩放保持范数，故
\[
 \|\mathbf1_Q\|_{M_p(\mathbb R^d)}=C_p^d,
 \quad1<p<\infty.
\]
这给出与矩形位置和各边长无关的界，但保留维数依赖。

为检查端点，取 \(\vartheta_j\in C_c^\infty(\mathbb R)\)，在 \([a_j,b_j]\) 的邻域恒为一，并令 \(\widehat f=\prod_j\vartheta_j\)。则 \(f\in\mathcal S\)，而其投影为各个区间逆变换的乘积。一维因子为
\[
 \mathcal F^{-1}\mathbf1_{[a,b]}(x)
 =e^{\pi\mathrm i(a+b)x}\cdot
       \frac{\sin\bigl(\pi(b-a)x\bigr)}{\pi x},
\]
原点按连续值补充。其绝对值不可积：每个充分远的周期中，正弦绝对值至少为 \(1/2\) 的部分占固定比例，相应积分下界形成发散的调和级数。Tonelli 因而使多维乘积的 \(L^1\) 范数也为无穷。这给出一个 \(L^1\cap L^2\) 输入，其 \(L^2\) 投影不属于 \(L^1\)，排除了 \(L^1\) 延拓。符号实值，规范 \(L^\infty\) 延拓若存在，其 \(L^1\) 预伴随正是同一符号的乘子，也将与上述结论矛盾。
\end{solution}

% 自拟；反例与尺度：全部消失矩也不消除同指数低频障碍，求带低频截断的逆算子范数增长。
\Needspace{6\baselineskip}
\begin{exercise}\label{b2:ex:22-new-5}
设整数 \(d\geq1\)、\(1<p<\infty\)、\(\alpha>0\)。对 Fourier 变换光滑紧支撑且远离原点的函数，定义
\[
 I_\alpha f=\mathcal F^{-1}\bigl[(2\pi|\xi|)^{-\alpha}\widehat f(\xi)\bigr].
\]
构造一族 \(f_r\in\mathcal S\)、\(r\downarrow0\)，使 \(\|f_r\|_p=1\)，每个多重指标 \(\beta\) 都满足
\(\int x^\beta f_r(x)\,\mathrm dx=0\)，但 \(\|I_\alpha f_r\|_p\to\infty\)。

\ProseParagraphBreak

再固定 \(\chi\in C_c^\infty(\mathbb R^d)\)，满足 \(0\leq\chi\leq1\)，在 \(|\xi|\leq1\) 为一、\(|\xi|\geq2\) 为零。对 \(\varepsilon>0\) 令
\[
 q_\varepsilon(\xi)=(2\pi|\xi|)^{-\alpha}
                    \cdot\bigl(1-\chi(\xi/\varepsilon)\bigr),
\]
在原点定义为零。证明其乘子范数恰为
\(\varepsilon^{-\alpha}\|q_1\|_{M_p}\)，其中 \(0<\|q_1\|_{M_p}<\infty\)。解释为何“有消失矩”与“具有固定正半径的频率缺口”是不同的控制。
\end{exercise}
\begin{solution}
取非零 \(\widehat f\in C_c^\infty(\{\,\xi:1<|\xi|<2\,\})\)，并归一化使 \(\|f\|_p=1\)。因为 \(\widehat f\) 在原点邻域恒零，其原点处的全部导数为零；Fourier 求导公式给出
\(\int x^\beta f(x)\,\mathrm dx=0\)。令
\[
 f_r(x)=r^{d/p}\cdot f(rx),\quad 0<r<1.
\]
缩放保持其 \(L^p\) 范数为一。每个 \(f_r\) 的 Fourier 变换仍在原点附近恒零，因而全部矩仍然消失。

令 \(g=I_\alpha f\)。频率乘积光滑紧支撑，所以 \(g\in\mathcal S\)，且非零。直接缩放给出
\[
 I_\alpha f_r(x)=r^{-\alpha}r^{d/p}\cdot g(rx),
 \quad\|I_\alpha f_r\|_p=r^{-\alpha}\|g\|_p\longrightarrow\infty.
\]
每一个输入都有频率缺口，但该缺口的半径与 \(r\) 一同趋零；消失矩本身没有给出统一的正半径。

符号 \(q_1\) 在原点附近为零，在固定过渡环带上光滑，在 \(|\xi|\geq2\) 上为 \((2\pi|\xi|)^{-\alpha}\)。其任意有限阶尺度导数都有界，故由 Mihlin 定理属于 \(M_p\)。它不是零符号，\cref{b2:prop:22-symbol-bound}保证其乘子范数为正。恒等式
\[
 q_\varepsilon(\xi)=\varepsilon^{-\alpha}
                                  \cdot q_1(\xi/\varepsilon)
\]
和符号缩放的等范数性质给出准确关系
\(\|q_\varepsilon\|_{M_p}=\varepsilon^{-\alpha}\|q_1\|_{M_p}\)。所以移除一个固定零频邻域确实给出同指数有界逆，但该邻域缩小时，常数具有明确且不可消去的增长。
\end{solution}

% 自拟；综合应用：Riesz变换组成向量场互补投影，核对符号、散度、旋度及平方空间的准确范数。
\Needspace{6\baselineskip}
\begin{exercise}\label{b2:ex:22-new-6}
设整数 \(d\geq2\)、\(1<p<\infty\)，向量空间 \(L^p(\mathbb R^d;\mathbb C^d)\) 使用逐点欧氏范数。令 \(R_j\) 为符号 \(-\mathrm i\xi_j/|\xi|\) 的 Riesz 变换，并对向量场 \(F=(F_1,\ldots,F_d)\) 定义
\[
 (QF)_j=-\sum_{k=1}^dR_jR_kF_k,
 \quad P=I-Q.
\]
证明 \(P,Q\) 有界，且
\[
 P^2=P,\quad Q^2=Q,\quad PQ=QP=0,\quad P+Q=I.
\]
在分布意义下验证
\[
 \sum_j\partial_j(PF)_j=0,
 \quad
 \partial_j(QF)_k-\partial_k(QF)_j=0.
\]
最后证明在 \(L^2\) 中两部分正交，两个投影的范数都等于一。可用紧支撑光滑标量函数的梯度和向量场 \((\partial_2\phi,-\partial_1\phi,0,\ldots,0)\) 检验两个下界。
\end{exercise}
\begin{solution}
Riesz 变换在 \(L^p\) 上有界，有限个分量的求和使 \(P,Q\) 在所给向量空间上有界，常数允许依赖 \(d,p\)。它们的矩阵符号分别为
\[
 q(\xi)=\frac{\xi\xi^{\mathsf T}}{|\xi|^2},
 \quad p(\xi)=I-\frac{\xi\xi^{\mathsf T}}{|\xi|^2},
 \quad\xi\ne0.
\]
这里 \((-\mathrm i)^2=-1\)，故定义中的负号使 \(q\) 具有正号。矩阵 \(q\) 是沿 \(\xi\) 方向的正交投影，\(p\) 是其正交补投影。因此上述四个代数恒等式先在 \(L^2\cap L^p\) 上由 Fourier 变换成立，再由稠密性推广到 \(L^p\)。

对 Schwartz 向量场，符号满足
\[
 \xi^{\mathsf T}p(\xi)=0,
 \quad \xi_jq_{k\ell}(\xi)-\xi_kq_{j\ell}(\xi)=0.
\]
乘上微分对应的 \(2\pi\mathrm i\) 后，这就是所需分布恒等式。对于一般 \(F\in L^p\)，取 Schwartz 逼近。投影的有界性给出输出的 \(L^p\) 收敛，而 \(L^p\) 收敛蕴含温和分布收敛；分布求导连续，所以两条恒等式仍成立。

在 \(L^2\) 中，逐频率正交分解与 Plancherel 给出
\[
 \|F\|_2^2=\|PF\|_2^2+\|QF\|_2^2,
 \quad\langle PF,QF\rangle_{L^2}=0.
\]
两投影因而均为压缩。取一个非恒定 \(\phi\in C_c^\infty\)，则非零梯度场 \(F=\nabla\phi\) 满足 \(QF=F\)，给出 \(\|Q\|=1\)。再选择 \(\phi\) 使 \((\partial_2\phi,-\partial_1\phi)\not\equiv0\)；该场的 Fourier 向量与 \(\xi\) 垂直，所以 \(PF=F\)，给出 \(\|P\|=1\)。在一般 \(L^p\) 中仍有有界的互补分解，但这里的正交恒等式专属于平方范数。
\end{solution}

% 自拟；中等：光滑频率截断与空间Lipschitz系数的交换子，明确初始定义并证明尖锐尺度。
\Needspace{6\baselineskip}
\begin{exercise}\label{b2:ex:22-new-7}
设整数 \(d\geq1\)，\(\chi\in C_c^\infty(\mathbb R^d)\)，\(k=\mathcal F^{-1}\chi\)，对 \(R>0\) 令
\(T_R=T_{\chi(\,\cdot\,/R)}\)。设 \(b:\mathbb R^d\to\mathbb C\) 为 Lipschitz 函数，常数为 \(L\)，不要求 \(b\) 有界。先对 \(f\in C_c^\infty\) 定义交换子
\[
 [b,T_R]f=b\cdot T_Rf-T_R(bf).
\]
证明它可由一个绝对积分定义于全部 \(L^p\)，并满足
\[
 \|[b,T_R]f\|_p
 \leq\frac L R\left(\int|y|\cdot|k(y)|\,\mathrm dy\right)\|f\|_p,
 \quad1\leq p\leq\infty.
\]
当 \(b(x)=x_j\) 且 \(\partial_j\chi\not\equiv0\) 时，求准确的交换子符号，证明对 \(1<p<\infty\) 其范数恰为一个与 \(R\) 无关的正数乘 \(R^{-1}\)。
\end{exercise}
\begin{solution}
核为 \(k_R(y)=R^d\cdot k(Ry)\)，它是 Schwartz 函数。在初始输入上，两项的差可写为
\[
 [b,T_R]f(x)
 =\int\bigl(b(x)-b(x-y)\bigr)\cdot k_R(y)\cdot f(x-y)\,\mathrm dy.
\]
因为 \(b\) 至多线性增长、\(f\) 紧支撑，初始两项有意义。差分满足
\(|b(x)-b(x-y)|\leq L|y|\)，从而绝对值由
\(L\bigl((|\,\cdot\,|\cdot|k_R|)*|f|\bigr)(x)\) 控制。这个非负卷积对有限指数输入几乎处处有限，对有界输入处处有限，因此上式也直接定义了一般输入上的交换子。

Young 不等式和换元给出
\[
 \|[b,T_R]f\|_p
 \leq L\left(\int|y|\cdot|k_R(y)|\,\mathrm dy\right)\|f\|_p
 =\frac L R\left(\int|z|\cdot|k(z)|\,\mathrm dz\right)\|f\|_p.
\]
有限指数时它与稠密延拓一致；无穷指数时使用上述直接积分，不使用 Schwartz 范数稠密性。若 \(b\) 无界，两个单独的算子项未必对全部 \(L^p\) 输入都有意义，真正延拓的是这个有抵消的差。

当 \(b(x)=x_j\) 时，核差分恰为 \(y_jk_R(y)\)。本卷 Fourier 约定下，
\[
 \mathcal F(y_jk_R)(\xi)
 =\frac{\mathrm i}{2\pi}\partial_{\xi_j}\bigl(\chi(\xi/R)\bigr)
 =\frac{\mathrm i}{2\pi R}\cdot(\partial_j\chi)(\xi/R).
\]
符号缩放保持乘子范数，所以
\[
 \|[x_j,T_R]\|_{p\to p}
 =\frac1{2\pi R}\|\partial_j\chi\|_{M_p}.
\]
导数符号光滑紧支撑，故范数有限；它非零，\cref{b2:prop:22-symbol-bound}使范数严格为正。因此一般的 \(R^{-1}\) 阶不能改善。
\end{solution}

% 自拟；中等：Bessel正核收缩的最优常数与不取到，区分每个输入严格和算子范数严格。
\Needspace{6\baselineskip}
\begin{exercise}\label{b2:ex:22-new-8}
设整数 \(d\geq1\)、\(s>0\)、\(1<p<\infty\)，\(J^{-s}\) 为符号 \((1+4\pi^2|\xi|^2)^{-s/2}\) 的 Bessel 位势算子。使用正文已建立的核 \(G_s\)，即
\(G_s>0\) 几乎处处、\(\int G_s=1\)、\(J^{-s}f=G_s*f\)，证明
\[
 \|J^{-s}\|_{p\to p}=1,
 \quad
 \|J^{-s}f\|_p<\|f\|_p
 \quad\text{若 }f\in L^p\setminus\{0\}.
\]
给出趋近算子范数的单位输入序列，并说明在 \(p=1\) 和 \(p=\infty\) 时为何不能保留上述逐个非零输入的严格不等式。
\end{exercise}
\begin{solution}
由正核的单位质量与 Jensen 不等式，几乎处处有
\[
 |G_s*f(x)|^p
 \leq\int G_s(y)\cdot|f(x-y)|^p\,\mathrm dy.
\]
Tonelli 给出右侧的总积分为 \(\|f\|_p^p\)，所以算子范数至多为一。若范数取等，两个非负可积函数的差积分为零，故上述 Jensen 不等式对几乎每个 \(x\) 都取等。选取一个等号成立且右侧有限的点 \(x\)。函数 \(z\mapsto|z|^p\) 在复数视作实二维向量时严格凸，因此取等迫使 \(f(x-y)\) 关于概率测度 \(G_s(y)\,\mathrm dy\) 几乎处处为同一个常数。因 \(G_s>0\) 几乎处处，\(f\) 也就关于 Lebesgue 测度几乎处处为常数。有限指数 \(L^p(\mathbb R^d)\) 中只有零常数符合要求，故每个非零输入都严格缩小。

为求算子范数下界，取 \(\psi\in C_c^\infty\)、\(\|\psi\|_p=1\)，并令
\[
 f_R(x)=R^{-d/p}\cdot\psi(x/R),\quad R\longrightarrow\infty.
\]
这些都是单位输入。缩放后用 Minkowski 不等式，得到
\[
 \|G_s*f_R-f_R\|_p
 \leq\int G_s(y)\cdot
                 \|\psi(\,\cdot-y/R)-\psi\|_p\,\mathrm dy.
\]
对每个固定 \(y\)，平移连续性使范数趋零；它又不超过 \(2\|\psi\|_p=2\)，而 \(G_s\) 可积。支配收敛给出误差趋零，所以
\(\|J^{-s}f_R\|_p\to1\)，算子范数恰为一。这同时说明范数并未由某个非零 \(L^p\) 输入达到。

当 \(p=1\) 时，对任意非零非负 \(f\in L^1\)，正性与 Tonelli 直接给出
\(\|G_s*f\|_1=\|f\|_1\)。当 \(p=\infty\) 时，常数输入满足 \(G_s*1=1\)。两个端点都有非零取等输入，不能把严格凸性与有限指数空间的论证直接搬过去。
\end{solution}
